% concatenated from ContinuousCoordination.tex
%% ****** Start of file apstemplate.tex ****** %
%%
%%
%%   This file is part of the APS files in the REVTeX 4.2 distribution.%%
%%   Copyright (c) 2024 The American Physical Society.
%%
%%   See the REVTeX 4 README file for restrictions and more information.
%%
% For Phys. Rev. appearance, change preprint to twocolumn.
% Choose physrev, prl, or rmp for journal
%  N.B. physrev is appropriate for all APS journals except prl and rmp
%  Add 'draft' option to mark overfull boxes with black boxes
%  Add 'showkeys' option to make keywords appear
\documentclass[aps,
preprint,
groupedaddress,
showkeys,
amsmath]{revtex4-2}
%\documentclass[aps,physrev,preprint,superscriptaddress]{revtex4-2}

\usepackage{graphicx}% Include figure files
\usepackage{esint}
\usepackage{tikz}
\usepackage{amsthm}
\usepackage{amssymb}
\usepackage{hyperref}

\newtheorem{theorem}{Theorem}
\newtheorem{theorem*}{Theorem}
\newtheorem{lemma}{Lemma}
\newtheorem{corollary}{Corollary}
\newtheorem{proposition}{Proposition}
\newtheorem{example}{Example}

\usetikzlibrary{shapes.geometric}

\DeclareMathOperator*{\argmax}{\text{argmax}}
\DeclareMathOperator*{\argmin}{\text{argmin}}
\DeclareMathOperator{\uu}{\mathbf{u}}
\DeclareMathOperator{\Rn}{\mathbb{R}^n}
\DeclareMathOperator*{\esssup}{ess\,sup}

% You should use BibTeX and apsrev.bst for references
% Choosing a journal automatically selects the correct APS
% BibTeX style file (bst file), so only uncomment the line
% below if necessary.
%\bibliographystyle{apsrev4-2}

\usepackage[normalem]{ulem}

\begin{document}

% Use the \preprint command to place your local institutional report
% number in the upper righthand corner of the title page in preprint mode.
% Multiple \preprint commands are allowed.
% Use the 'preprintnumbers' class option to override journal defaults
% to display numbers if necessary
%\preprint{}

%Title of paper
\title{Nonlinear Nonlocal Diffusion Equations for the Analysis of Continuous Coordination and Anti-Coordination Type Games}

% repeat the \author .. \affiliation  etc. as needed
% \email, \thanks, \homepage, \altaffiliation all apply to the current
% author. Explanatory text should go in the []'s, actual e-mail
% address or url should go in the {}'s for \email and \homepage.
% Please use the appropriate macro foreach each type of information

% \affiliation command applies to all authors since the last
% \affiliation command. The \affiliation command should follow the
% other information
% \affiliation can be followed by \email, \homepage, \thanks as well.
\author{John S. McAlister}
\email{Contact author: jmcalis6@vols.utk.edu}
\altaffiliation[Also at ]{the National Institute for Modeling Biological Systems}
\author{Nina H. Fefferman}
\altaffiliation[Also at ]{the National Institute for Modeling Biological Systems}
\altaffiliation[ ]{the Center for Analysis and Prediction of Pandemic Expansion (APPEX)}
\altaffiliation[ and the ]{Department of Ecology and Evolutionary Biology, University of Tennessee - Knoxville}
\author{Tadele A. Mengesha}
\affiliation{Department of Mathematics, University of Tennessee - Knoxville}

%Collaboration name if desired (requires use of superscriptaddress
%option in \documentclass). \noaffiliation is required (may also be
%used with the \author command).
%\collaboration can be followed by \email, \homepage, \thanks as well.
%\collaboration{}
%\noaffiliation

\date{\today}

\begin{abstract}
Coordination games with explicit spatial or relational structure are of interest to economists, ecologists, sociologists, and others studying emergent global properties in collective behavior. When assemblies of individuals seek to coordinate action with one another through myopic best response or other replicator dynamics, the resulting dynamical system can exhibit many rich behaviors. However, these behaviors have been studied only in the case where the number of players is countable and the relational structure is described discretely. By giving an extension of a general class of coordination-like games, including true coordination games themselves,  into a continuous setting, we can begin to study coordination and cooperative behavior with a new host of tools from PDEs and nonlocal equations. In this study, we propose a rigorously supported extension of structured coordination-type games into a setting with continuous space and continuous strategies and show that, under certain hypotheses, the dynamics of these games are described through a nonlinear, nonlocal diffusion equation. We go on to prove existence and uniqueness for the initial value problem in the case where no boundary data are prescribed. For true coordination games, we go further and prove a maximum principle, weak regularity results, as well as some numerical results toward understanding how solutions to the coordination equation behave. We present several modeling results, characterizing stationary solutions both rigorously and through numerical experiments and conclude with a result towards the inhomogeneous problem. 
\end{abstract}

\keywords{Coordination, Nonlinear, Nonlocal diffusion}

\maketitle

\section{Introduction}\label{SEC:Introduction}
In this paper we will discuss a broad class of games, which we call Toeplitz games, where payoffs for each pairwise interaction depend only on the  ``distance" between the strategies of each player. Neutral Coordination games (games wherein players receive a payoff which is maximized when they play the same strategy as their co-player but is independent of the strategy itself) fall under this category, as well as neutral anti-coordination and dis-coordination games.  Each of these games, but especially the coordination game, has been of particular interest since the rise of evolutionary game theory because of the rich behaviors observed when this game is considered dynamically. 
Early results considering coordination games show that, in well-mixed populations, the consensus equilibrium (wherein every individual plays the same strategy) is the only long-term stable equilibrium \cite{Kandori1993}. Adding spatial structure by way of considering players as vertices of a graph returns similar results in the case where one strategy may dominate the other, either in terms of payoff or risk \cite{Ellison1993, Ellison2000}. A general review of coordination with spatial structure before 2010 can be found in \cite{Weidenholzer}. Moreover, there has been numerical work on spatially explicit coordination in larger, more general graphs \cite{Buskens2016, Raducha2022, McAlister2024}. When examining the critical case where every option is neutral, there are many non-consensus equilibria. However, even in this setting, such equilibria become exceedingly rare to discover for very large graphs.

All of these previous studies, like the ones mentioned above as well as \cite{Arditti2024,Berninghaus2010,Clauset2004, Ely2002,Gilboa1991,Oechssler1997,Oechssler1999,Robson1995}, have considered discrete strategies with discrete player spaces. These models are helpful for thinking of particular examples, (e.g. the selection of a coordinated computer operating system among coworkers in an office) but are unable to be translated to conditions where strategy may vary continuously (e.g. language conventions). Further, using a continuous player space to represent general behaviors of replaceable individuals organized in space, rather than discrete player spaces with explicit interaction behavior unique to each player, changes the required analysis for the system. The purpose of the present study is to introduce a rigorous continuous extension of Toeplitz-type games so that the questions in the application areas may be investigated with a new set of tools.

When other games have previously been considered with continuous strategies, integral equations have been of use. Indeed, many of these authors have extended beyond the present study to consider mixed continuous strategies in the space of distributions (i.e. Stackelberg equilibrium, nonlocal replicator dynamics \cite{Kavallaris2018}) but these often rely on a finite number of players. In the same way, studies of a continuum of players often require discrete strategy. In the present study we seek to understand a system with continuous players and strategies that cannot be considered as a potential game \cite{SANDHOLM2001, Quang2016}. Moreover, The non-monotonicity and nonlinearity of the nonlocal equation in our problem is of interest beyond its applications here.  

In particular, we propose a way to translate the Toeplitz-type games, into continuous strategic and player domains through the use of a nonlinear, nonlocal diffusion equation. We start in section \ref{SEC:Extension} with some game theoretic background, then we present a rigorously supported extension into the continuous setting resulting in an nonlocal equation model. In section \ref{SEC:ExistenceUniquenessAndMainResults} we prove classical existence and uniqueness of solutions. We strengthen these results for true coordination games in subsection \ref{SUBSEC:CoordinationToeplitzGames} by way of a weak maximum principle. In addition to this, in subsection \ref{SUBSEC:AdditionalResults}, we consider the problem in the Cauchy setting to show strengthened results. In particular, we get regularity estimates so that, after showing several analytical examples in section \ref{SEC:AnalyticalExamples}, we can approximate them through simple numerical methods in section \ref{SEC:NumericalResults} and consider other properties of the nonlocal equations. In section \ref{SEC:ModelingResults}, we consider what the theory tells us about coordination in continuous space by examining stationary solutions and conducting several numerical examples. Finally, in section \ref{SEC:Inhomogeneous} we note that this model is compatible with an inhomogeneity which provides a way for us to extend the model to an even larger class of games. 

\section{Continuous Extension}\label{SEC:Extension}
\subsection{Game Theoretic Background}\label{SUBSEC:GameTheoreticBackground}
In a strategic form game, there are a set of players, $V$, which each have a set of strategies, $S_v$. A strategy profile is a collection of players' strategies from the Cartesian product $S:=\prod_{v\in V}S_v$. Each player also has a payoff function $w_v:S\rightarrow \mathbb{R}$. For our investigation, we are interested in games where every player has access to the same set of strategies $S$ and each player has the same payoff function $w:S^{|V|}\rightarrow \mathbb{R}$. 

Crucial to the study of strategic form games is the concept of\textit{ best response}. For a focal individual $v$ and a strategy profile $s\in S$, we break the strategy profile down as $(s_v,s_{-v})$ which is the ordered pair of player $v$'s strategy and the strategies of the players other than $v$ respectively. We say that $v$'s \textit{best response} to $s$ is whatever strategy in $S_v$ maximizes $v$'s payoff supposing all the other players play according to $s_{-v}$. That is

\[ br_v(S)=\argmax_{c\in S_v}\{w_v(c,s_{-v})\}.\]

If players play mixed strategies, which are probability distributions over a set of pure strategies, then we make the distinction between best response $br_v(s)$, which is the (possibly mixed) strategy that maximizes payoff, and pure strategy best response $BR_v(s)$, which is the pure strategy which maximizes payoff. A last crucial definition which will be obvious to a game theory audience is that $s\in S$ is called a Nash equilibrium if and only if all players are playing a best response to $s$. That is $s$ is a Nash equilibrium $\iff s=(s_1,s_2,...,s_n)$ and $s_v\in br_v(s)$ for all $v$.

We are interested in a class of games which we call Toeplitz games. These are the discrete two player games which have Toeplitz payoff matrices. Equivalently each player has a set of pure strategies which can be ordered as $S_1=S_2=\{s^{(1)},s^{(2)},...,s^{(m)}\}$ so that if the payoff  matrix $A$ is given where $a_{i,j}$ describes the payoff $w(s^{(i)},s^{(j)})$ against strategy $j$, $A$ is diagonal-constant (AKA Toeplitz). Three examples below from left to right are the pure coordination game, an anti-coordination game, and a game which is neither coordination nor anti-coordination. 

\[\begin{bmatrix}
    1&0&0&0\\0&1&0&0\\0&0&1&0\\0&0&0&1
\end{bmatrix}
\quad\quad\quad
\begin{bmatrix}
    0&1&2&4\\
    1&0&1&2\\
    2&1&0&1\\
    4&2&1&0
\end{bmatrix}
\quad\quad\quad 
\begin{bmatrix}
    0&1&2&3\\
    -1&0&1&2\\
    -2&-1&0&1\\
    -3&-2&-1&0
\end{bmatrix}\]
The ordering required to write the payoff matrix as a Toeplitz matrix, $(s^{(1)},s^{(2)},...,s^{(m)})$ induces a metric on the strategy space $d(s^{(i)},s^{(j)})=|i-j|$. With a metric and a complete ordering, this kind of game can be reduced to a two player game in which each player $v$ selects a number $n_v\in\{1,...,m\}$ and their payoff is determined by $w_v(n_v,n_u)=\rho(n_v-n_u)$. This ordering is not unique. for instance, in the pure coordination game, any ordering of the strategies will satisfy this condition. 

Those Toeplitz games which are coordination games are of particular interest to us. A two-player coordination game is a type of strategic form game that satisfies the Bandwagon Property\cite{Kandori1998}. This property is when both players use the same strategy set $S$, for each player, and for any (possibly mixed) strategy $s\in S$, $BR_v(s)\subseteq C(s)$, where $C(s)$ is the support of $s$. For a pure strategy, the support is just the strategy itself $C(s)=\{s\}$. For a mixed strategy, the support is all of those pure strategies which are expressed with strictly positive probability. Similarly an anti-coordination game is a two-player game wherein, if the payoff function is made negative, the resulting game satisfies the bandwagon property. 

These coordination and anti-coordination games which can be expressed as a Toeplitz game are of particular interest because they model the homogeneous impact of cooperative or uncooperative behavior. In a Teoplitz game, a particular strategy does not give any fitness benefit independent of its proximity to a co-player's strategy. That is to say that any fitness benefit is only the result of the interaction between players (whether the interactions are sympathetic or antagonistic). If the strategies were reorganized by some permutation $p$ so that all the players playing $s^{(l)}$ take on the strategy $s^{(p(l))}$, the payoff for every player will have remained the same. If several methods of communication are equally efficient but the benefit of communication is only achieved when two players are using the same method, this can be described as a Toeplitz coordination game. Building a foundational theory of this homogeneous problem will then allow for better treatment of inhomogeneous coordination and anti-coordination processes.

\subsection{Continuous Extension in Space}\label{SUBSEC:ContinuousExtensionsInSpace}
As described above, dyadic interactions of this type are easy to understand. Therefore, when we seek to understand multiplayer Toeplitz games, we easily generalize the two-player interactions and say that a player's payoff is the sum (or equivalently the arithmetic mean) of the payoffs from each dyadic interaction.
As is typical, we will start by considering the players as vertices in a graph where (possibly weighted) edges describe the strength of interaction between individuals. Suppose the game is played on the graph $G(V,E)$ with weighted adjacency matrix $W$ and that each player has pure strategies $B=\{\hat{e}_1,\hat{e}_2,...,\hat{e}_m\}$, which is the standard basis for $\mathbb{R}^m$. Then, if we consider a strategy profile as a function $u:V\rightarrow B$ rather than an element of the Cartesian product $B^{|v|}$ (although the two spaces are clearly identical), we can write our payoff function as 
\begin{equation}\label{disPayoff}
    w(v|u)=\sum_{i\in V}W_{i,v}u(v)^TAu(i)    
\end{equation}
where $A$ is the payoff matrix as described above. This formulation holds true for any symmetric game with a payoff matrix $A$, it need not be Toeplitz. 

If instead of a collection of discrete players we considered an uncountable continuum of players, we might think of our player domain as a subset of $\mathbb{R}^n$, and instead of summing we will integrate over the entire domain. The weighted adjacency matrix is replaced by a nonnegative integrable kernel $K\in C_b^0(\Omega; L^1(\mathbb{R}^n))$ (That is a kernel $K(x,y)$ which is continuous $x$, in the $L^1$ sense, and integrable in $y$ where $\sup_{x\in\Omega}\|K(x,\cdot)\|_{L^1(\mathbb{R}^n)}<\infty$). Replacing the sum in \eqref{disPayoff} with the integral we get that a strategy profile $u:\Omega\rightarrow B$ gives the payoff
\begin{equation}\label{ctsSpacePayoff}
w(x|u)=\int_\Omega K(x,y)u(x)^TAu(y)dy.
\end{equation}
Again, this extension does not require the game to be a Toeplitz game. 

\subsection{Continuous Strategic Extensions}\label{SUBSEC:ContinuousExtensionsInStrategy}
To truly study cooperative and non-cooperative behavior we may also consider extending the existing model to include continuous strategies. There are two ways to do this, one which does not require the use of a Toeplitz type game (which we will mention but not discuss in depth) and the other which is exceptionally helpful when we have a Teoplitz type game. 

The first is the \textit{mixed strategy} concept. As before, if players take on a strategy which is a probability distribution over the pure strategies, then we can use the same payoff function \eqref{ctsSpacePayoff} but allow for $u$ to map from $\Omega $ to $\Delta^{m-1}:=\{x\in[0,1]^m;\sum_{i=1}^mx_i=1\}.$ This is a natural way to think of continuity in the strategy space and is worthy of further study, but it is not the focus of the present paper.

The second is what we call the \textit{comparable strategy} concept. If our payoff matrix is a Toeplitz matrix, we can express the pairwise payoff simply as some function $\rho(dd(u(x),u(y))$ where $dd$ is the directed metric which is natural from the ordering induced by the Teoplitz matrix.  The continuous extension from this point is clear. The ordering of the strategies means that they can be considered as elements of $\mathbb{R}$ and we can allow $\rho:\mathbb{R}\to \mathbb{R}$ so that our payoff can be expressed as $\rho(u(x)-u(y))$.  Therefore we can write our payoff function as 
\begin{equation}\label{ctsStratPayoff}
w(x|u)=\int_{\Omega}K(x,y)\rho(u(x)-u(y))dy.
\end{equation}
To make this extension all we have done is replace the bilinear form in \eqref{ctsSpacePayoff} with our payoff function induced by the Teoplitz matrix. We are especially interested in this type of game, both because of its implications in the application areas but also because, in the continuous form, we will see that it has great similarities to a class of important nonlocal equations called nonlocal diffusive equations. Indeed the time dependent model will be a nonlinear, nonlocal diffusion equation and understanding this model will contribute to the understanding of nonlocal diffusion problems in general.  

With this fitness function, we have now described a strategic form game completely as we have a set of players $\Omega$, a set of strategies $\mathbb{R}$, and a payoff function \eqref{ctsStratPayoff}. However, the search for Nash equilibria to this game is exceedingly difficult and the primary difficulty is the size of the function space in which we must work. Notice that if $\rho(z)=\chi_{\{0\}}(z)$ then any constant function is a Nash equilibrium to this game. This is clear because if $u\equiv c$ then 
\begin{equation*}
\begin{split} BR_x(u)&=\argmax_{z\in\mathbb{R}}\{\int_\Omega K(x,y)\chi_{\{0\}}(z-c)dy\}\\ 
    &=\argmax_{z\in\mathbb{R}}\{\chi_{\{0\}}(z-c)\|K(x,\cdot)\|_{L^1(\Omega)}\}\\
    &=c
    \end{split}
    \end{equation*}
    so $BR_x(u)=u(x)$ in $\Omega$. However, if $\tilde{u}$ differs from $u$ as a non-empty set of measure 0 then it will certainly not be a Nash equilibrium. This causes a problem because it means that we cannot search for equilibria in any $L^p$ space (as elements of these spaces are equivalence classes of functions that differ at sets of measure 0). 

    Because of this issue, in order to understand the game in this setting we will follow in the example of those who have studied the game in the discrete setting and consider it as a dynamic game through myopic best response (e.g \cite{Ellison1993,Raducha2022}). Myopic best response is a replication dynamic for an evolutionary game in which a set of players, who chosen to update their strategies, take on their best response (often pure strategy best response) to the current strategy profile.  By repeating this process we can study a time series of strategy profiles which, if it terminates, will result in a strategy profile in which every player is playing a best response. Under myopic best response we can see that the evolution of strategy profiles for the game with payoff function \eqref{ctsStratPayoff} will evolve according to a particular nonlocal equation.

    \begin{proposition}\label{PROP:ModelWellFounded}
        Under myopic best response, bounded strategy profiles of the game with players $\Omega\subseteq \mathbb{R}^n$ choosing strategies in $\mathbb{R}$, with fitness as in \eqref{ctsStratPayoff}, will evolve as
       
       \[\frac{\partial}{\partial t}u(x,t)=\int_{\Omega}K(x,y)\rho'(u(x,t)-u(y,t))dy\]
        so long as the following three hypotheses are met
        \begin{itemize}
            \item[(H1)] Players change their strategies in arbitrarily small time steps $\Delta t$
            \item[(H2)] Players incur a quadratic cost $\frac{h^2}{\Delta t}$ for changing their strategy
            \item[(H3)] $\rho\in C^{1,1}(\mathbb{R})$ and $\rho(z)\leq Cz^2+A$ for some nonnegative $C,A$
        \end{itemize}
    \end{proposition}

    \begin{proof}
        Consider a bounded strategy profile $u(\cdot, t):\Omega\to \mathbb{R}$. Every player will seek to update their strategy by some amount $h$ in order to take on their best response to $u(\cdot, t)$ after a time step of $\Delta t$ and in doing so they will incur a cost of $\frac{h^2}{2}$. Let 
        
       \[S_x(h):=\int_\Omega K(x,y)\rho(u(x,t)+h-u(y,t))dy\] be the payoff that player $x$ will receive after updating their strategy by $h$. Because $\rho$ is subquadratic (H3), we can see that 

        \begin{equation*}
        \begin{split}
            S_x(h)&\leq C\int_\Omega K(x,y)(h+(u(x,t)-u(y,t))^2dy+\underbrace{A\sup_{x\in\Omega}\|K(x,\cdot)\|_{L^1(\Omega)}}_{A_1}\\
            &\leq \underbrace{C\int_\Omega K(x,y)dy}_{C_1}h^2 + 2hC\int_\Omega K(x,y)(u(x,t)-u(y,t))dy\\
            &\quad +C\int_\Omega K(x,y)(u(x,t)-u(y,t))^2dy+A_1
        \end{split}
        \end{equation*}
        Since $K\in C_b^0(\Omega; L^1(\mathbb{R}^n))$, $C_1$ and $A_1$ can be bounded by constant independent of $x$. Now because $u$ is bounded we know that $|u(x,t)-u(y,t)|$ is bounded by some $M$ and thus we can write 
        \begin{equation*}
            \begin{split}
                S_x(h)&\leq C_1h^2+\underbrace{2CM\sup_{x\in\Omega}\|K(x,\cdot)\|_{L^1(\Omega)}}_{B_1}h+ \underbrace{CM^2\sup_{x\in\Omega}\|K(x,\cdot )\|_{L^1(\Omega)} + A_1}_{A_2}\\
                &\leq C_1h^2+B_1h+A_2
            \end{split}
        \end{equation*}
        Having shown that $S_x(h)$ is uniformly subquadratic in $h$, we also note that when $\rho\in C^{1,1}$ (H3), $S_x(h)$ is also in $C^{1,1}$. Observe that 
        \begin{equation*}
            \frac{d}{dh}S_x(h)=\int_\Omega K(x,y)\rho'(u(x,t)+h-u(y,t))dy
        \end{equation*}
        And so clearly for any compact subdomain $I\subset \mathbb{R}$ we can write 
        \begin{equation}\label{EQ:SinC11}
        \begin{split}
            \bigg|&\frac{d}{dh}S_x(h_1)-\frac{d}{dh}S_x(h_2)\bigg|\\
            &\leq \int_\Omega |K(x,y)| |\rho'(u(x,t)+h_1-u(y,t))-\rho'(u(x,t)+h_2-u(y,t))|dy\\
            &\leq \int_\Omega K(x,y)L_{\rho}|h_1-h_2|dy\\
            &\leq \sup_{x\in\Omega}\|K(x,\cdot)\|_{L^1(\Omega)}L_{\rho}|h_1-h_2|
            \end{split} 
        \end{equation}
        Because $u$ is bounded, the choice of a compact subdomain $I$ gives us a Lipschitz constant for $\rho'$ on the subdomain $[-2\sup |u|-\sup I,2\sup|u|+\sup I]$ which is called $L_\rho$ in \eqref{EQ:SinC11}. Therefore $\frac{d}{dh}S_x(h)$ is locally Lipschitz. 

        This is important because it means that, for any $x$, when $\Delta t$ is small enough, $S_x(h)-\frac{h^2}{\Delta t}\leq C_2h^2+A_3$ for some negative $C_2$ and some $A_3$. Thus $S_x(h)-\frac{h^2}{\Delta t}$ must have a global maximizer, $h^*$ and that global maximizer will satisfy 
        
       \[\frac{d}{dh}S_x(h^*)=2\frac{h^*}{\Delta t}\]
        because everything is continuously differentiable. 

        Having observed this, we note also that there must be a negative $h^-$ so that when $h<h^-$ then $S_x(h)-\frac{(h)^2}{\Delta t}<S_x(0)$. Likewise there must be an $h^+$ so that $h>h^+\implies S_x(h)-\frac{(h)^2}{\Delta t}<S_x(0)$.
        These $h^-$ and $h^+$ will form a compact interval which will surely contain $h^*$. Moreover, from the chosen interval we have a Lipschitz constant for $\frac{d}{dh}S_x$ which we call $L_S$.
        The next step is to put bounds on $h^*$, in order to do this we will consider different cases: when $h^*>0, h^*<0$, and $h^*=0$. 

        If $h^*>0$ we write that 
        \begin{equation*}
            \begin{split}
                -h^*L_S&\leq \frac{d}{dh}S_x(h^*)-\frac{d}{dh}S_x(0)\leq h^*L_S\\
                -h^*L_S&\leq \frac{h^*}{2}-\frac{d}{dh}S_x(0)\leq h^*L_S\\
                \end{split}
        \end{equation*} 
        Reorganization on each side of the inequality, with the knowledge that $\Delta t$ can be made small enough so that $2-L_S\Delta t>0$, will yield 
        \begin{equation*}
                \frac{d}{dh}S_x(0)\frac{\Delta t}{2+L_S\Delta t}\leq h^*\leq \frac{d}{dh}S_x(0)\frac{\Delta t}{2-L_S\Delta t}
        \end{equation*}

        Likewise we can show that if $h^*<0$ we will get the inequality 
        \begin{equation*}
            \frac{d}{dh}S_x(0)\frac{\Delta t}{2-L_S\Delta t}\leq h^*\leq \frac{d}{dh}S_x(0)\frac{\Delta t}{2+L_S\Delta t}
        \end{equation*}

        We also note here that if $h^*=0$ this requires that $\frac{d}{dh}S_x(0)=0$. 

        Now that we have bounds on $h^*$, recall that $h^*$ is the change in strategy in one time step. That it $h^*=u(x,t+\Delta t)-u(x,t)$. If we make this substitution and divide our inequalities by $\Delta t$ we see that 
        \begin{equation*}
            \begin{split}
                \frac{d}{dh}S_x(0)\frac{1}{2+L_S\Delta t}&\leq \frac{u(x,t+\Delta t)-u(x,t)}{\Delta t}\leq\frac{d}{dh}S_x(0)\frac{1}{2-L_S\Delta t}\quad \quad h^*>0\\
                \frac{d}{dh}S_x(0)\frac{1}{2-L_S\Delta t}&\leq \frac{u(x,t+\Delta t)-u(x,t)}{\Delta t}\leq\frac{d}{dh}S_x(0)\frac{1}{2+L_S\Delta t}\quad \quad h^*<0
            \end{split}
        \end{equation*}
        Both inequalities trivially hold in the case that $h^*=0$. In any case, when we take $\Delta t\to 0$ we see that by squeeze theorem 
        \begin{equation*}
            \frac{\partial}{\partial t}u(x,t)=\frac{1}{2}\frac{d}{dh}S_x(0)
        \end{equation*}

        To complete the proof we need only note that 
        
       \[\frac{d}{dh}S_x(0)=\int_\Omega K(x,y)\rho'(u(x,t)-u(y,t))dy\] and do a trivial rescaling of space-time to arrive at the desired nonlocal equation.
        \end{proof}

        Thus we have shown that, for any Toeplitz type game, we can make a continuous extension and describe how bounded strategy profiles will evolve in time with a nonlocal equation. We will call the nonlocality 
        \begin{equation}\label{nonlocality}
            g[u](x,t):=\int_\Omega K(x,y)\rho'(u(x,t)-u(y,t))dy
        \end{equation}
        and express the nonlocal equation as $u_t=g[u]$. Notice that if $K$ and $\rho$ are selected appropriately a certain Toeplitz coordination game extends directly to the nonlocal heat equation \cite{NonlocalDiffusionProblems,GOMEZ2017} and a certain Toeplitz anti-coordination game extends directly to the nonlocal backward heat equation. Being able to study these games with a new suite of tools from PDEs and nonlocal equations will allow us to understand coordination and anti-coordination in space more fully. Moreover, the present investigation into this game will also extend our understanding on nonlinear nonlocal diffusion equations.  
\section{Existence, Uniqueness, and Other Main Results for Toeplitz Type Games}\label{SEC:ExistenceUniquenessAndMainResults}
\subsection{General Toeplitz Games}\label{SUBSEC:GeneralToeplitzGames}
 Now that we have a nonlocal equation which captures the behavior we are interested in, we will proceed with some classical existence and uniqueness results for the initial value problem (IVP) wherein an initial strategy profile $u_0$ is given in $\Omega$ and there are no boundary data prescribed. The assumption is that no information is entering the system from outside the domain and so it is argued in \cite{CHASSEIGNE2006} that this is the analogue to the Neumann boundary condition for nonlocal diffusion type problems. It is important to note that these results for the nonlocal equation do not depend on the assumptions that make the model appropriate for the application area ($H1-3$ from proposition \ref{PROP:ModelWellFounded}). Namely, $\rho$ need not be subquadratic. However, we will require that $\rho\in C^{1,1}$. 
 
 The main result of this section is a Picard iteration type proof relying on the contraction mapping principle. In order to achieve this, we will first state and prove two important lemmas. The first will show that the nonlocality, $g$, maps continuous and  bounded functions to continuous and bounded functions. The second will show that $g$ is locally Lipschitz with respect to the sup norm in $C^0_b$. For the following, let $\Omega\subset \mathbb{R}^n$ be an open domain and let $\Omega_T:=\Omega\times [0,T]$. 

  \begin{lemma}[$g$ is well defined]\label{LEM:welldefined}
        The non local operator $g[u](x,t)=\int_\Omega K(x,y)\rho'(u(x)-u(y))dy$ is well defined from $C^0_b(\Omega_T,\mathbb{R})$ into $C^0_b(\Omega_T,\mathbb{R})$.
    \end{lemma}

    \begin{proof}
        To show that $g$ is well defined in this space, let $u\in C^0_b(\Omega_T)$ and we will show that $g[u]$ is continuous on $\Omega_T$ and bounded. To show that $g[u]$ is bounded, first note that, when $u$ is bounded, because $\rho\in C^{1,1}$, surely there is a $W$ such that $|\rho'(u(x,t_1)-u(y,t_2))|<W$ for all $x,y\in \Omega$ and $t_1,t_2\in [0,T]$ by a continuous on compact argument. This means that $|g[u](x,t)|\leq \int_\Omega |K(x,y)|Wdy\leq W\|K(x,\cdot)\|_{L^1(\Omega)}$. By assumption $\|K(x,\cdot)\|_{L^1(\Omega)}$  is uniformly bounded so $\sup_{x\in\Omega}\|K(x,\cdot)\|_{L^1(\Omega)}\leq D$ for some $D<\infty$. Thus we have shown that $\|g[u]\|_\infty<WD$, so $g[u]$ is bounded whenever $u$ is continuous and bounded.

        Consider the sequence ${(x_n,t_n)}_{n=1}^\infty\subset\Omega_T$ with $(x_n,t_n)\rightarrow (x,t)\in\Omega$. Furthermore, consider the difference 

        \begin{equation*}
        \begin{split}
        |g[u](x_n,t_n)&-g[u](x,t)|\\
        &=\bigg|\int_\Omega K(x_n,y)\rho'(u(x_n,t_n)-u(y,t_n))-K(x,y)\rho'(u(x,t)-u(y,t))dy\bigg|\\
            &\leq \int_\Omega|(K(x_n,y)-K(x,y))\rho'(u(x_n,t_n)-u(y,t_n))|dy\\
            &\quad +\int_\Omega |K(x,y)(\rho'(u(x_n,t_n)-u(y,t_n))-\rho'(u(x,t)-u(y,t)))|dy\\
            &\leq I_1+I_2
            \end{split}
        \end{equation*}
        We will consider each integral separately. Let $\epsilon>0$. Note that there is an $R$ such that $\int_{\Omega\setminus B_{R}(0)}K(x,y)dy<\frac{\epsilon}{8W}$. Now consider $x_n$ and $y$ inside $\overline{B_R(0)}.$ $\rho\in C^{1,1}$ so $\exists \delta$ such that $|x-y|<\delta \Rightarrow |\rho'(x)-\rho'(y)|<\frac{\epsilon}{4D}$. Because $u$ is continuous, $\exists M_1$ such that $|u(x_n,t_n)-u(x,t)|<\delta/2$ whenever $n>M_1$. Furthermore, In $\Omega\cap \overline{B_R(0)}$ there is surely a $M_2$  such that $n>M_2\implies |u(y,t_n)-u(y,t)|<\delta/2$ for any $y\in \Omega\cap\overline{B_R(0)}$. It is an easy application of the triangle inequality to see that this implies $|(u(x_n,t_n)-u(y,t_n))-(u(x,t)-u(y,t))|<\delta$ when $n>\max\{M_1,M_2\}$ and thus, when $n$ is sufficiently large
        
       \[|\rho'(u(x_n,t_n)-u(y,t_n))-\rho'(u(x,t)-u(y,t))|\leq \frac{\epsilon}{4D} \quad\forall y\in\Omega\cap \overline{B_R(0)}.\]

        Therefore we can see that
        \begin{equation*}
            \begin{split}
                I_2&\leq 2W\int_{\Omega\setminus B_R(0)}K(x,y)dy+\int_{\Omega\cap B_R(0)}K(x,y)\frac{\epsilon}{4D}dy\\
                &\leq 2W\frac{\epsilon}{8W}+D\frac{\epsilon}{4D}=\frac{\epsilon}{2}
            \end{split}
        \end{equation*}

        Now we consider $I_1$. It is clear that
        
       \[I_1 \leq W\int_{\Omega}|K(x_n,y)-K(x,y)|dy.\] 
        From the assumed continuity of $\|K(x,\cdot)\|_{L^1(\Omega)}$ we get immediately that $\exists M_3$ such that for $n>M_3$, 
        $\|K(x_n,\cdot)- K(x,\cdot)\|_{L^1(\Omega)}\leq \frac{\epsilon}{2W}$.

        Thus, if $M=\max\{M_1,M_2,M_3\}$, then $n>M\implies$
        
       \[|g[u](x_n,t_n)-g[u](x,t)|\leq I_1+I_2\leq \frac{\epsilon}{2}+W\frac{\epsilon}{2W}=\epsilon\]
        This clearly works for any $\epsilon>0$ so we have shown that $g[u]$ is continuous at $(x,t)$ for any $(x,t)\in \Omega _T$. Therefore we have proved that $g:C^0_b(\Omega_T,\mathbb{R})\rightarrow C^0_b(\Omega_T,\mathbb{R})$ is well defined.  
    \end{proof}
    
    \begin{lemma}[Lipschitz Continuity of $g$] \label{LEM:Lipschitzg}
        For any bounded (in the sup norm sense) subset $X_R:=\{u\in C^0_b(\Omega_T);\|u\|_\infty \leq R\}\subset C^0_b(\Omega_T)$, there exists a $C^g\geq 0 $ such that for every $u,v\in X_R$ 
        \begin{equation*}
            \|g[u](\cdot,t)-g[v](\cdot,t)\|_{\infty} \leq C^g\|u(\cdot,t)-v(\cdot,t)\|_{\infty} 
        \end{equation*}
    \end{lemma}

    \begin{proof}
        Let $u,v\in X_R$. Then we observe that for each $0<t<T,$
        \begin{equation*}
            \begin{split}
              \|g[u](\cdot,t)-g[v](\cdot,t)\|_\infty   &=\bigg\|\int_\Omega K(x,y)(\rho'(u(x,t)-u(y,t))-\rho'(v(x,t)-v(y,t))dy\bigg\|_\infty\\
                &\leq \sup_{x\in \Omega}\int_\Omega|K(x,y)||\rho'(u(x,t)-u(y,t))-\rho'(v(x,t)-v(y,t))|dy\\
                &\leq \sup_{x\in \Omega}\|\rho'(u(x,
                t)-u(\cdot,t))-\rho'(v(x,t)-v(\cdot,t))\|_\infty \int_\Omega K(x,y)dy. 
            \end{split}
        \end{equation*}
        
        
        Recall that $\rho\in C^{1,1}(\mathbb{R})$  and $\rho'$ has Lipschitz constant $L_\rho$ for the compact interval $[-2R,2R]$. Now note that for any $(x,t)\in\Omega_T$
        \begin{equation*}
            \begin{split}
                \sup_{y\in \Omega}|\rho'(u(x,t)&-u(y,t))-\rho'(v(x,t)-v(y,t))|\\
                &\leq \sup_{y\in \Omega}L_\rho|u(x,t)-u(y,t)-v(x,t)+v(y,t)|\\
                &\leq \sup_{y\in \Omega}L_\rho (|u(x,t)-v(x,t)|+|u(y,t)-v(y,t)|)\\
                &\leq L_\rho(|u(x,t)-v(x,t)|+\|u(\cdot,t)-v(\cdot, t)\|_\infty)
            \end{split}
        \end{equation*}
        so naturally 
        
       \[\sup_{x\in \Omega}\|\rho'(u(x,t)-u(\cdot,t))-\rho'(v(x,t)-v(\cdot,t))\|_\infty \leq 2L_\rho\|u(\cdot,t)-v(\cdot,t)\|_{\infty}\]
        This, and the fact that $\|K(x,\cdot)\|_{L^1(\Omega)}\leq D$ uniformly for some finite $D$, gives us the result that for any $t\in[0,T]$
        
       \[\|g[u](\cdot,t)-g[v](\cdot,t)\|_{\infty}\leq 2L_\rho D\|u(\cdot,t)-v(\cdot,t)\|_{\infty}\]
        Clearly, then we have $C^g=2L_\rho D$ and 
        
       \[\sup_{t\in [0,T]}\|g[u](\cdot,t)-g[v](\cdot,t)\|_{\infty}\leq C^g\sup_{t\in [0,T]}\|u(\cdot,t)-v(\cdot,t)\|_{\infty}.\] It is important to notice that $L_\rho$ may depend on the choice of $R$ so $C^g$ depends on $R$.
        \end{proof}

        Having shown that $g$ is well defined and Lipschitz continuous we can now prove short time existence and uniqueness of the solution to the initial value problem $u_t=g[u]$ in $\Omega_T$ with $u(0,t)=u_0\in C^0_b(\Omega)$ through a contraction mapping principle.


    \begin{theorem}[Short Time Existence and Uniqueness] \label{THM:ShortTimeExistenceAndUniqueness}
        The initial value problem $u_t=g[u]$ in $\Omega_\tau$ has a unique continuous and bounded solution in $\Omega_\tau$ for some $\tau$, when $u(x,0)=u_0\in C_b^0(\Omega)$, $\rho\in C^{1,1}(\mathbb{R})$, and $K\in C^0_b(\Omega;L^1(\Omega))$.
    \end{theorem}
    \begin{proof}
        Let $\Omega_T=\Omega\times [0,T]$ with $T$ to be chosen later. Equip the function space $C_b^0(\Omega_T)$ with the standard sup norm $\|u\| = \sup_{t\in[0,T]}\|u(\cdot,t)\|_\infty$. 
        Now, for some $R>\|u_0\|_\infty$, let $E_{R,T}:=\{u\in C^0_b(\Omega_T,\mathbb{R});u(x,0)=u_0, \|u\|\leq R\}$. Observe that $E_{R,T}$ is nonempty as the map $(x,t)\mapsto u_0(x)$ belongs in  $E_{R,T}$. Moreover, observe that $E_{R,T}$ is complete with respect to the sup norm so we will be able to proceed with a Banach Fixed Point Theorem (BFPT) argument. It is clear that a solution to the IVP will also satisfy
        
       \[u(x,t)=u_0(x)+\int_0^tg[u](x,s)ds.\]
        Let $\Theta:C_b^0(\Omega_T,\mathbb{R})\rightarrow C_b^0(\Omega_T,\mathbb{R})$, where $\Theta u = u_0+\int_0^tg[u]ds$, and notice that if this operator has a unique fixed point in $E_{R,T}$ then we have a unique solution to the IVP. 

        First, we show that $\Theta:E_{R,T}\rightarrow E_{R,T}$ for some $T$. It is easy to say that $\Theta u\in C^0_b(\Omega_T,\mathbb{R})$ and that $\Theta u(x,0)=u_0(x)$.  Because of lemma \ref{LEM:welldefined} we know that $g[u]\in C_b^0(\Omega_T)$ so its time antiderivative is obviously continuous in space and time. 
        
        In order to show that that $\|\Theta u\| \leq R$ we note that $g[u]$ is bounded whenever $u$ is bounded. In particular, $g[u]\leq WD$ where $W$ is the bound for $\rho'$ on  $[-2R,2R]$ which, of course, depends on $R$, and $D$ is the uniform bound on $\|K(x,\cdot)\|_{L^1(\Omega)}$. Therefore we can always find a $T_R$ such that $\|\Theta u(\cdot, t)\|_\infty \leq \|u_0\|_\infty +\int_0^tWD \leq R$ for all $t\in [0,T_R]$. Let $T<T_R$ and we have that $\Theta:E_{R,T}\rightarrow E_{R,T}$. 
        
        Now that we have shown that $\Theta$ indeed maps from $E_{R,T}$ to $E_{R,T}$ for some $T$, we need only show that $\Theta$ is a contraction in $E_{R,T}$ for some, possibly smaller, value of $T$. Let $u,v \in E_{R,T}$ and note that
        \begin{equation*}
        \begin{split} 
            \|\Theta u-\Theta v\| &= \sup_{t\in[0,T]} \bigg\|\int_0^tg[u](\cdot,s)-g[v](\cdot,s)ds\bigg\|_{\infty}\\
            &\leq \sup_{t\in[0,T]}\int_0^t\|g[u]-g[v]\|ds
            \end{split} 
        \end{equation*}
        By lemma \ref{LEM:Lipschitzg} we get immediately that
        
       \[\|\Theta u-\Theta v\|\leq C^gT\|u-v\|.\] Notice here that $C^g$ depends on $R$ because the Lipschitz constant for $g$ is defined for a particular compact subset.  
        With this Lipschitz constant for $\Theta$, when $T\leq\frac{1}{2C^g}$ we know that $\Theta$ is a contraction. Thus, by the BFPT  there is a unique $u\in E_{R,T}$ such that $\Theta u = u$.  Thus, there is a continuous and bounded $u$ so that $u_t=g[u]$ on $\Omega_T$ and $u(0,t)=u_0\in C^0_b(\Omega)$ so long as $T<\min\{T_R,\frac{1}{2C^g}$\}. This completes the proof of short time existence and uniqueness by way of the Banach Fixed Point Theorem. We can use an extension principle to get longer existence time. Note that for a given $u_0$ we can choose an $R_1$ and find a resulting $T_1$ so that there is a unique solution on $[0,T_1]$. Take $u(\cdot, T_1-\epsilon)$ for some $\epsilon>0$ as our initial condition, take a new $R_2$ and resulting $T_2$ to find a new solution on $[T_1-\epsilon, T_2]$. Wherever these solutions overlap they must be identical because of the uniqueness proved here. We cannot use this extension to say that the solution is global in time because we have no lower bound on the minimal existence time. As $R$ grows, the resulting $T$ may shrink quickly enough so that we cannot extend the solution beyond some finite time $\tau$.     
    \end{proof}

    We have shown that for any continuously extended Teoplitz game with certain hypotheses on $K$ and $\rho$, and with a bounded and continuous initial strategy profile, there is exactly one way the strategy profile will evolve for some time. Next we will show that if there is a finite maximal existence time, $T$, (i.e. $u(x,t)$ cannot be extended beyond time $T$ as a solution of the IVP) then there is necessarily finite time blowup.

    \begin{lemma}[Finite Time Blow up]\label{LEM:FiniteTimeBlowup}
    If  $T<\infty$ is the maximal time of existence for a solution $u$ to the IVP $u_t=g[u]$ with $u(\cdot,0)=u_0\in C^0_b(\Omega)$ then $\|u(\cdot, t)\|_\infty \rightarrow \infty $ as $t\to T$.
    \end{lemma}

    \begin{proof}
        Suppose that $u$ is a solution to the IVP $u_t=g[u]$ with $u(\cdot, 0)=u_0\in C^0_b(\Omega)$ that has a maximal time of existence $T<\infty$. Moreover, by way of contradiction suppose there is a bounded subset $E_R:=\{v\in C^0_b(\Omega);\|v\|_\infty\leq R\}$ such that for all $t\in [0,T)$ $u(\cdot, t)\in E_R$. Note that $E_{R}$ is  a closed subset of $C^0_b(\Omega).$

        Consider now a sequence of times $t_n$ which have $t_n\to T$ as $n\rightarrow \infty$. This is necessarily a Cauchy sequence in the reals. We will show that $u(\cdot, t_n)$ is a Cauchy sequence in $E_R$ with respect to the sup norm. Note that for any $\delta>0$ there is an $N$ so that $n,m>N\implies |t_n-t_m|<\delta$. 
        \begin{equation*}
            \begin{split}
                \|u(\cdot,t_n)-u(\cdot, t_m)\|_\infty&=\bigg\|\int_0^{t_n}g[u](\cdot, s)ds-\int_0^{t_m}g[u](\cdot,s)ds\bigg\|_\infty\\
                &=\bigg\|\int_{t_n}^{t_m}g[u](\cdot,s)ds\bigg\|_\infty\\
                &\leq \delta \sup_{s\in[t_n,t_m]}\|g[u](\cdot,s)\|_\infty
            \end{split}
        \end{equation*}
        By assumption, $\|u\|_\infty$ is bounded for all time $t\in [t_n,t_m]\subset [0,T)$ uniformly in time by $R$. By lemma \ref{LEM:welldefined} we know that if $u$ is bounded by $R$ then there is a $R_g$ so that $\|g[u]\|_\infty\leq R_g$ and, crucially this upper bound only depends on the choice of $R$ (see proof of lemma \ref{LEM:welldefined}).  Thus we can say that $\sup_{s\in[t_n,t_m]}\|g[u](\cdot,s)\|_\infty\leq R_g$. Thus we have that
        
       \[\|u(\cdot, t_n)-u(\cdot, t_m)\|_\infty\leq \delta R_g\]

        Therefore, for any $\epsilon>0$, we let $\delta=\frac{\epsilon}{2R_g}$ and from this we get an appropriate $N$ so that $|t_n-t_m|<\delta$ and thus $\|u(\cdot,t_n)-u(\cdot, t_m)\|\leq \epsilon$. Thus $u(\cdot, t_n)$ is a Cauchy sequence from the closed subset of a Banach space, $E_R$, and so it has a limit $\tilde{u}_0\in E_R$. (Note carefully that $\tilde{u}_0$ is continuous because the convergence is uniform). Now consider another arbitrary sequence of times $\tau_k$ with $\lim \tau_k\rightarrow T$ and we will show that $u(\cdot, \tau_k)\rightarrow \tilde{u}_0$. Of course, $u(\cdot,t)$ is differentiable in time on $[0,T)$ and that derivative is uniformly bounded by $R_g$ as above, so $|u(x,t)-u(x,\tau)|\leq R_g|t-\tau|$. Thus we can write 
        \begin{equation*}
            \begin{split}
                \|u(\cdot, \tau_k)-\tilde{u}_0\|&\leq \|u(\cdot, \tau_k)-u(\cdot, t_n)\|_\infty+\|u(\cdot,t_n)-\tilde{u}_0\|_\infty\\
                &\leq R_g|\tau_k-t_n|+\|u(\cdot,t_n)-\tilde{u}_0\|_\infty
            \end{split}
        \end{equation*}
        Both $\tau_k$ and $t_n$ approach $T$ from the left and so if $\epsilon>0$ there is an $n_1$ such that $k,n>n_1\implies |t_k-t_n|<\frac{\epsilon}{2R_g}$. Furthermore, we already showed that $\|u(\cdot,t_n)-\tilde{u}_0\|\rightarrow 0$ so there is an $n_2$ such that $ n>n_2\implies \|u(\cdot,t_n)-\tilde{u}_0\|_\infty \leq \frac{\epsilon}{2}$. Therefore we get that when $k>\max\{n_1,n_2\}$ then surely
        
       \[\|u(\cdot, \tau_k)-\tilde{u}_0\|\leq \epsilon.\]
        This is true for any sequence of $\tau\rightarrow T$ from the left thus we write that $\lim_{t\to T^-}u(\cdot,t)=\tilde{u}_0\in C^0_b(\Omega)$.
        
        
        Because we have a $\tilde{u}_0\in C^0_b(\Omega)$, by theorem \ref{THM:ShortTimeExistenceAndUniqueness} we can find a solution $\tilde u$ to the IVP $\tilde u=g[\tilde u]$ with $\tilde u(\cdot, T)=\tilde u _0$ on some interval $[T,T+\eta).$ Now let $\hat u(\cdot, t) = u(\cdot,t)$ when $t\in [0,T)$ and $\hat u(\cdot, t)=\tilde u(\cdot ,t)$ when $t\in [T,T+\eta)$. We will show now that $\hat u\in C^0_b(\Omega_{T+\eta}).$ The boundedness is immediate from the hypothesis and from the details of the proof of theorem \ref{THM:ShortTimeExistenceAndUniqueness}. Also, we know that $\hat{u}$ is continuous on $[0,T)$ and on $(T,T+\eta)$. To see that is also continuous at $T$ observe that $\lim_{t\rightarrow T^-}\hat{u}(\cdot, t)=\hat{u}(\cdot, T)$ by the previous result and that $\lim_{t\rightarrow T^+}\hat{u}(\cdot, t)=\hat{u}(\cdot, T)$ by the details in the proof of \ref{THM:ShortTimeExistenceAndUniqueness}. This convergence is uniform and so we can say that  $\hat u\in C^0_b(\Omega_{T+\eta}).$
        
        Lastly, we show that $\hat{u}$ is a solution to the IVP, that is, in the integral form, 
        
       \begin{equation}\label{integralequation}\hat{u}(x,t)=u_0(x)+\int_0^tg[\hat{u}](x,s)ds.\end{equation}
        This is obviously true when $t<T$. Moreover we can see that when $t=T$ we can write
        \begin{equation}\label{corEq2}
            \begin{split}
            u_0(x)+\int_0^Tg[\hat{u}](x,s)ds&=u_0(x)+\int_0^{T-\epsilon}g[{u}](x,s)ds+\int_{T-\epsilon}^Tg[{u}](x,s)ds\\
                &=u(x,T-\epsilon)+\int_{T-\epsilon}^Tg[u](x,s)ds
            \end{split}
        \end{equation}
        Recall that, because $\|\hat{u}\|_\infty=\|u\|_\infty\leq R$ when $t<T$,  $\|g[\hat{u}]\|_\infty=\|g[u]\|_\infty\leq R_g$ when $t<T$. Therefore if we take the limit as $\epsilon \rightarrow 0$ in equation \eqref{corEq2} we get $u_0(x)+\int_0^Tg[u](x,s)ds=\lim_{t\rightarrow T^-}(x,t)=\tilde{u}_0(x)=\hat{u}(T)$. 

        Finally, when $t>T$ we observe that 
        \begin{equation*}
            \begin{split}
                u_0(x)+\int_0^tg[\hat{u}](x,s)ds&= u_0(x)+\int_0^Tg[\hat{u}](x,s)ds+\int_T^tg[\hat{u}](x,s)ds\\
                &=\tilde{u}_0+\int_T^tg[\tilde{u}](x,s)ds\\
                &=\tilde{u}(x,t)=\hat{u}(x,t)
            \end{split}
        \end{equation*}

        Therefore $\hat{u}\in C^0_b(\Omega_{T+\eta})$ and $\hat{u}$ satisfies that integral equation \eqref{integralequation}. Thus $\hat{u}$ extends the solution to the IVP $u_t=g[u]$ with $u(\cdot,0)=u_0$ beyond the assumed maximal time of existence. This contradicts the assumption that $T$ was the maximal time of existence. Thus, if there is a finite maximal time of existence, the solution must leave every compact subset of $C^0_b(\Omega_T)$. Because the limit $u(\cdot,t_n)$ is continuous for $t_n\leq T<\infty$ we know that $\|u(\cdot, t)\|_\infty\to \infty$ as $t\to T$.
    \end{proof}
    Note that this result only discusses the consequences of having a finite time of existence. It is made less interesting by the fact that we have no example of a solution which exhibits finite time blowup. 
    These results are for general continuous extensions of Toeplitz games with appropriate $\rho$. These general games are interesting but we can achieve more specific results when we constrain our study to Toeplitz coordination games.  


\subsection{Coordination Toeplitz Games}\label{SUBSEC:CoordinationToeplitzGames}
In the discrete case, if a Toeplitz game has a positive Strictly Diagonally Dominant (SDD) payoff matrix then it is necessarily a coordination game. We can see this because, for any mixed strategy $t\in \Delta^{m-1}$, the payoff of playing a pure strategy $k\in \{1,...,m\}$ is $w(e_k,t)=e_k^TAt$. It is immediate to see that if $k$ is such that $t_k=0$ then ${e_k^{T}At}=\sum_{j=1}^ma_{kj}t_j<a_{kk}\|t\|_{\ell^\infty}$ and if $l$ is such that $t_l=\|t\|_{\ell^\infty}$ then ${e_l^{T}At}=\sum_{j=1}^ma_{lj}t_j\geq a_{ll}t_l=a_{kk}\|t\|_{\ell^\infty}$. Therefore if $k$ is not in the support of a strategy $t$ it cannot be a pure strategy best response to $t$ and so $BR(t)\subset C(t)$. 

Although the translation between an SDD matrix and the function $\rho$ used to describe pairwise payoff is not clear, we consider the continuous coordination Toeplitz game as games where
\begin{equation}\label{coordinationcondition}
\rho'(z)\begin{cases}\leq0&z>0\\
=0&z=0\\
\geq 0 & z<0\end{cases}
\end{equation} 
This implies that $\rho$ achieves its global maximum at $z=0$ and so the pairwise interaction with pure strategies governed by such a $\rho$ will satisfy the bandwagon property (and thus, this is indeed a coordination game). This restriction on $\rho$ allows us to think more specifically about how coordination behavior evolves in time. 

It is appropriate to make the comparison between the continuous Toeplitz coordination games and nonlocal diffusion equations, which are well covered in \cite{NonlocalDiffusionProblems}. Indeed, if we choose $\rho(z)=\frac{-1}{2}z^2$ our coordination equation becomes exactly the linear nonlocal diffusion equation $u_t=\int_\Omega K(x,y)(u(y)-u(x))dy$. Moreover, regardless of our choice of $\rho$, if we assume it is even and $C^{1,1}$, near zero $\rho'(z)\approx-z.$ Because of this similarity, we will expect some of the same behavior as the nonlocal diffusion equation. In particular, we will see that there is a weak maximum principle, and thus solutions to the IVP will exist globally in forward time. In the same way, it is appropriate to compare a continuous Toeplitz anti-coordination game to the backward nonlocal diffusion equation.

\begin{lemma}[Weak Maximum Principle]\label{LEM:MaximumPrinciple}
        If $u$ solves $u_t = g[u]$ in $\Omega_T$ with $u(\cdot, 0)=u_0\in C^0_b(\Omega)$ and if $\rho$ satisfies \eqref{coordinationcondition}, then 
       \[\|u(\cdot, t_2)\|_\infty\leq\| u(\cdot , t_1)\|_\infty \] whenever $t_1\leq t_2$. 
    \end{lemma}
    \begin{proof}
        Notice that it is no loss of generality to assume $t_1=0$, so we will prove that $\|u(\cdot,t)\|_\infty\leq \|u_0\|_\infty$ for all $t\in [0,T)$ where $T$ is the maximal existence time which may be infinite. Let $\epsilon>0$ and let $v=\|u_0\|_\infty +\epsilon$. Now observe that if $u$ solves $u_t=g[u]$ and $\tilde{u}:=u-\epsilon t$, then $\tilde{u}_t=g[\tilde{u}]-\epsilon$. This is because $g[u-\epsilon t]=g[u]$ (More generally, $g$ is invariant under vertical shifts, even if they are time-dependent). Suppose that $\tilde{u}(x^\star,t^\star) = v$ for the first time at some $x^\star\in \Omega$ and some $t^\star>0$. That means that when  $t<t^\star$, $\tilde{u}(y,t)<v$ for all $y\in \Omega$ and by continuity of $u$, $\tilde{u}(y,t^\star)\leq v$ for all $y\in\Omega$. Using the fact that $\rho'(z)\leq 0$ when $z\geq 0$, we obtain that    
        \[g[\tilde{u}](x^\star,t^\star)=\int_\Omega K(x,y)\rho'(\tilde u(x^\star,t^\star)-\tilde u(y,t^\star))dy\leq 0, \] 
        and therefore, $\partial_t\tilde u(x^\star,t^\star)\leq 0-\epsilon<0$. However, because $\tilde{u}$ is continuously differentiable in time (see proof of theorem \ref{THM:ShortTimeExistenceAndUniqueness}) and because $\tilde u=v$ for the first time at $t=t^\star$ we can say $\partial_t\tilde u(x^\star,t^\star)\geq 0$. This is a contradiction so we can say surely that $\tilde{u} < v$ which means $u< \|u_0\|_\infty+\epsilon+\epsilon t$. This is true for any $\epsilon$ so let $\epsilon\rightarrow 0$ and we see that $\|u(\cdot, t)\|_\infty<\|u_0\|_\infty$ for all finite time. This proves the result.
    \end{proof}

    This is tremendously helpful because it will allow us to give a uniform lower bound on existence time if we seek to extend solutions forward in time. Because we know that, if we select an $R\in \mathbb{R}$, a solution to the initial value problem with $\rho$ satisfying \eqref{coordinationcondition} will exist on a time interval $[0,T)$ and surely have $\|u(\cdot,T-\epsilon)\|_\infty<R$ and so it will be in the same closed subset of $C^0_b(\Omega_T)$, $E_{R,T}=\{u\in C^0_b(\Omega_T);u(\cdot,t)=u_0, \|u\|_\infty\leq R\}$. Because the provable existence time depended only on the forms of $K$, $\rho$, and the bound $R$, when we consider a coordination game, we can extend by the same amount of time in each iteration. We can repeat the process indefinitely to get global existence. Global existence can be proven this way, but we may also prove the same result just by considering the previous two lemmas.

    \begin{theorem}[Global existence and uniqueness with particular $\rho\in C^{1,1}$]\label{THM:GlobalExistenceAndUniqueness}
    Let $\rho\in C^{1,1}(\mathbb{R})$ satisfy \eqref{coordinationcondition}. Under this strengthened hypothesis, the Initial Value Problem $u_t=g[u]$ with $u(x,0)=u_0\in C_b^0(\Omega)$ has a unique continuous and bounded solution for all finite time.
    \end{theorem}

    \begin{proof}
        Observe that for any time $t$, the solution $u(\cdot t)$ must be in the compact subset 
        \[E_0:=\{u\in C_b^0(\Omega);\|u\|_\infty\leq \|u_0\|_\infty\}\subset C^0_b(\Omega)\]
        because of lemma \ref{LEM:MaximumPrinciple}. If the maximal time of existence is $T<\infty$ then it must leave this compact subset before time $T$ by lemma \ref{LEM:FiniteTimeBlowup}. This is a contradiction so $u$ cannot have a finite maximal time of existence. 
    \end{proof}

    Lemma \ref{LEM:MaximumPrinciple} and Theorem \ref{THM:GlobalExistenceAndUniqueness} are consistent with our understanding of the coordination game in the discrete case. For multiplayer coordination, it is an easy extension of the bandwagon property to say that innovation outside of the support of the current strategy profile is never a best response. Indeed, lemma \ref{LEM:MaximumPrinciple} is the continuous version of the Weak Bandwagon Property of \cite{cui2022}, which says that it is never optimal for an individual to take on a strategy not used by any opponent in a multiplayer coordination game. 

    The boundedness of solutions for the coordination game enables us to give a weak regularity result.  This result will be especially important when we attempt to approximate solutions though numerical methods. 


    \begin{theorem}[Regularity for the coordination game without boundary conditions]\label{THM:Regularity}
        Let $\Omega\subseteq \mathbb{R}^n$. Suppose $u$ solves the Initial Value Problem $u_t=g[u]$ with the coordination assumption \eqref{coordinationcondition} and with $u_0\in C^{0,1}_b(\Omega)$ with a uniform Lipschitz constant $L_0$, If $K$ is uniformly Lipschitz in the first variable with respect to the $L^1$ norm, then for any finite time $t$, $u(\cdot, t)\in C^{0,1}(\Omega)$ . Moreover, if $\Omega$ is open and bounded with a $C^1$ boundary,  $u\in W^{1,\infty}(\Omega)$. Furthermore, for some positive $c$ and $C$, 
        
       \[\|D_xu(\cdot, t)\|_{L^\infty(\Omega)}\leq (L_0+Ct)e^{ct} \]
        Where $C$ depends only on $K$ and $u_0$ and $c$ depends only on $K$ and $\rho$.
    \end{theorem}

    \begin{proof}
        Let $\phi(h,x,t)=u(x+h,t)-u(x,t)$ for some $h\in  \mathbb{R}^n$ so that both $x$ and $x+h$ are in $\Omega$  and observe that by simply subtracting the two solutions from one another.
        \begin{equation}\label{regproofeq1}
            \phi(h,x,t)=\phi(h,x,0)+\int_0^t\left(g[u](x+h,s)-g[u](x,s)\right)ds
        \end{equation}
        Let $L_\rho$ be the Lipschitz constant for $\rho'$ on the interval $[-2\|u_0\|_\infty,2\|u_0\|_\infty]$. Recall by lemma \ref{LEM:MaximumPrinciple} that $u\in [-\|u_0\|_\infty,\|u_0\|_\infty]$ for all $x,t\in \Omega_T$ so surely $u(x,t)-u(y,t)\in [-2\|u_0\|_\infty,2\|u_0\|_\infty]$ regardless of $x$ and $y$. Also note that $\rho'$ is bounded by some $W$ on the same interval. Now observe that this integrand is
        \begin{equation*}
            \begin{split}
                |g[u]&(x+h,s)-g[u](x,s)|\\
                &=\bigg|\int_\Omega K(x+h,y)\rho'(u(x+h,s)-u(y,s))-K(x,y)\rho'(u(x,s)-u(y,s))dy\bigg|\\
                &\leq \int_\Omega K(x+h,y)|\rho'(u(x+h,s)-u(y,s))-\rho'(u(x,s)-u(y,s))|dy\\
                &\quad \quad +\int_\Omega [K(x+h,y)-K(x,y)]|\rho'(u(x,s)-u(y,s))|ds\\
                &\leq \int_\Omega K(x+h,y)L_\rho|u(x+h,s)-u(x,s)|dy\\
                &\quad\quad+W\|K(x+h,\cdot)-K(x,\cdot)\|_{L^1(\Omega)}\\
                &\leq L_\rho\int_\Omega K(x+h,y)|\phi(x,h,s)|dy+WL_Kh\\
                & \leq L_\rho\sup_{z\in \Omega}\|K(z,\cdot)\|_{L^1(\Omega)}\|\phi(h,\cdot,s)\|_{\infty}+WL_kh\\
            \end{split}
        \end{equation*}
        let $c=L_\rho\sup_{z\in \Omega}\|K(z,\cdot)\|_{L^1(\Omega)}$ and $C=WL_k$ and make this replacement into equation \eqref{regproofeq1}  and take the sup norm over $\Omega$ to see that
        
       \[\|\phi(h,\cdot,t)\|_\infty \leq \|\phi(h,\cdot,0)\|_\infty +\int_0^t[c\|\phi(h,\cdot, s)\|_\infty +Ch]ds\]
        
        It is an easy application of Gr\"onwall's inequality to see that
        \begin{equation*}
             \|\phi(h,\cdot,t)\|_\infty
             \leq(\|\phi(h,\cdot, 0)\|_\infty+Cht)e^{ct}
        \end{equation*}
        Note also that $\|\phi(h,\cdot,0)\|_\infty\leq |h|L_0$, so the difference from any $x\in \Omega$ on any finite time interval $[0,T]$ is controlled by
        
        \[
            |u(x,t)-u(x+h,t)|\leq |h|(L_0+Ct)e^{cT}
        \]

        Notice that we have shown that for each $t$, $u(\cdot,t)$ is globally Lipschitz continuous in $\Omega$ with Lipschitz constant $L_T=(L_0+CT)e^{cT}$. Note that this constant is uniform in $t$. This means, when $\Omega$ is open and bounded, we can use the characterization of $W^{1,\infty}(\Omega)$ \cite{evansPDE} to conclude that $u$ is weakly differentiable and the weak derivative is bounded by the Lipschitz constant.
        
         \[\|D_xu(\cdot,t)\|_{L^\infty(\Omega)}\leq L_T=(L_0+Ct)e^{cT}\]
         
    \end{proof}

    Note that, in the above proof, the Lipschitz continuity of $u$ holds without restrictions on the domain. Indeed, if $\Omega$ is not bounded we have a global Lipschitz constant. Also note that in any open domain $\Omega$ and for each $t,$ $u(\cdot, t)\in W^{1,\infty}_{loc}(\Omega)$.

    \begin{corollary}
        For $\Omega \subseteq \mathbb{R}^n$, suppose that $u$ solves the IVP $u_t=g[u]$ under the assumptions in theorem \ref{THM:Regularity} on a finite time domain $[0,T]$. In this case, $u$ is globally Lipschitz in $\Omega\times[0,T]$
    \end{corollary}

    \begin{proof}
        By theorem \ref{THM:Regularity}, $u$ is Lipschitz continuous in space and the Lipschitz constant is uniform on a bounded time interval. Moreover, we know that $u(x,\cdot)$ is continuously differentiable because $u_t(x,t)=g[u](x,t)$ which is continuous in space and time by \ref{LEM:welldefined}. $g[u]$ is bounded above by a constant related to the form of $K,\rho$ and the bounds on $u$. Under the hypothesis of theorem \ref{THM:Regularity}, $\|u(\cdot, t)\|_\infty\leq \|u_0\|_\infty$ so we obtain  
        \begin{equation*}
            \begin{split}
                |u(x,t)-u(y,s)|&\leq |u(x,t)-u(y,t)|+|u(y,t)-u(y,s)|\\
                &\leq C_T|x-y|+C_{0}|t-s|\\
                &\leq L(|x-y|+|t-s|)
            \end{split}
        \end{equation*}
        Where $L$ is a constant which depends only on $T, K, \rho$ and $u_0$. Of course we can adjust this norm to say that 
        \[|u(x,t)-u(y,s)|\leq C \sqrt{|x-y|^2+|t-s|^2}\]
        so that we can express it using the euclidean norm in $\Omega_T\subset \mathbb{R}^{n+1}$
    \end{proof}

    \subsection{The Cauchy Problem with a translation invarient kernel}\label{SUBSEC:AdditionalResults}
    The results presented in the previous subsections are rather general and do not rely on heavy assumptions about the domain or the form of $K$ or $\rho$ beyond what is strictly necessary for the model to be well posed. The analysis of this model is made exceedingly difficult, however, by the nonlinear nature of the nonlocality. Because the nonlocality is both nonlinear and non-monotonic we currently have no comparison principle as in \cite{Santos2022} nor can we use Fourier analysis or semigroup theory as is the standard for linear nonlocal diffusion problems \cite{NonlocalDiffuisionandApplications,NonlocalDiffusionProblems,Kavallaris2018} to directly analyze the model. To take our analysis further we present a strengthening of the regularity in the Cauchy setting with a different assumption on $K$.
  
    Much of the existing literature on nonlocal problems focuses on the use of translation invariant or even radial kernels. In that tradition we will examine the improved result we can obtain through using a kernel of the form $K(z)\in L^{1}(\Omega)$. This assumption means that every player has the same pattern of interaction distributed spatially. In the case that $\Omega=\Rn$ we can mildly strengthen the regularity result from theorem \ref{THM:Regularity}. The purpose of this strengthening is to remove the dependence on the shape of $K$ from the Lipschitz constant. In future work, we will seek to investigate the ``zero-horizon limit" or the ``local limit" of this nonlinear nonlocal diffusion problem (i.e., the limit as the support of $K$ goes to $\{0\}$). In the case that $\rho'$ is linear, we see that the results on non-local limits from \cite{DU2015,Du2022,Mengesha2015,NonlocalDiffusionProblems,Tao2017,Tao2019} will hold. For the nonlinear case, the scaling of the kernel to achieve the non-local limit, poses a problem for the provable regularity of solutions which has, at present, prevented us from characterizing the local limit of this diffusion equation. However,  this strengthened regularity result will be crucial in this pursuit.  

    \begin{theorem}[Regularity for the coordination game with a Translation Invariant Kernel]\label{THM:CauchyRegularity}
    Consider the domain $\Rn$. Suppose that $u$ solves the Initial Value Problem $u_t=g[u]$ with the coordination assumption \eqref{coordinationcondition} and with $u_0=C_b^{0,1}(\Rn)$ with the uniform Lipschitz constant $L_0$. If $K$ is a translation invariant Kernel (i.e. $K(x,y)=J(x-y)$) then $u(\cdot, t)\in C^{0,1}(\Rn)$ for any finite time. Moreover the global Lipschitz constant for $u(\cdot, t)$ on $ \Rn$ is given as $L_0e^{cT}$ where $c$ depends only on the Lipschitz constant for $\rho'$ and on $\|u_0\|_\infty$. 
    \end{theorem}

    \begin{proof}
        Without loss of generality, we normalize the kernel $J(z)$ so that $\|J\|_{L^{1}(\Rn)}=1$.
        We will proceed with this proof in much the same way as theorem \ref{THM:Regularity}. As before we let $\phi(h,x,t) = u(x+h,t)-u(x,t)$ for some $h\in \mathbb{R}^n$ with some magnitude $r$ and a bearing $\theta \in S^{n-1}$ and again observe the equation \eqref{regproofeq1} holds. 

        Let $L_\rho$ be the Lipschitz constant for $\rho'$ on the interval $[-2\|u_0\|_\infty, 2\|u_0\|_\infty]$ and by lemma \ref{LEM:MaximumPrinciple} we know that $u\in [\|u_0\|_\infty, \|u_0\|_\infty]$ for all $x,t\in \Omega_T$. We take the same computation as before to find that the quantity 
        \begin{equation*}
            \begin{split}
               &|g[u](x+h,s)-g[u](x,s)|\\ \leq &\underbrace{\int_{\Rn} J(x+h-y)|\rho'(u(x+h,s)-u(y,s))-\rho'(u(x,s)-u(y,s))|dy}_{I_1}\\
                & \quad +\underbrace{\int_{\Rn} |(J(x+h-y)-J(x-y))\rho'(u(x,s)-u(y,s))|dy}_{I_2} 
            \end{split}
        \end{equation*}
        As before, $I_1\leq L_\rho\|J(x-\cdot)\|_{L^1(\Rn)}\|\phi(h,\cdot,s)\|_\infty =L_\rho\|\phi(h,\cdot, s)\|_\infty$. For $I_2$ we can see that
        \begin{equation*}
            \begin{split}
                I_2&=\bigg|\int_{\Rn} [J(x+h-y)-J(x-y)]\rho'(u(x,s)-u(y,s))dy\bigg|\\
                &=\bigg|\int_{\Rn} J(x+h-y)\rho'(u(x,s)-u(y,s))dy-\int_{\Rn} J(x-y)\rho'(u(x,s)-u(y,s))dy\bigg|\\
                &=\bigg|\int_{\Rn+h} J(x-y)\rho'(u(x,s)-u(y-h,s))dy-\int_{\Rn} J(x-y)\rho'(u(x,s)-u(y,s))dy\bigg|\\
                &=\bigg|\int_{\Rn} J(x-y)[\rho'(u(x,s)-u(y-h,s))-\rho'(u(x,s)-u(y,s))]dy\bigg|\\
                &\leq \int_{\Rn} J(x-y)L_\rho\big|u(x,s)-u(y-h,s)-u(x,s)+u(y,s)|dy\\
                &\leq L_\rho\|J\|_{L^1(\Rn)}\|u(\cdot -h,s)-u(\cdot, s)\|_{L^\infty(\Rn)}\\
                &\leq L_\rho \|\phi(-h,\cdot,s)\|_{L^\infty(\Rn)}
            \end{split}
        \end{equation*}

    Now instead of considering only the supremum of $\phi(-h,x,s)$ for $x\in \Rn$, we will decompose $h$ into $r$ and $\theta$ and consider 
    
   \[\psi(r,s):=\sup_{x,\theta\in \Rn\times S^{n-1}}\phi(r,\theta,x,s).\]
    From our inequalities on $I_1$ and $I_2$ We can see that 
    
   \[\phi(h,\cdot, t)\leq \phi(h,\cdot, 0)+\int_0^tL_\rho\|\phi(h,\cdot,s)\|_{L^\infty(\Omega)}ds+\int_0^tL\|\phi(-h,\cdot, s)\|_{L^\infty(\Omega)}ds\]
    So we can make the replacement that 
    
   \[\psi(r,t)\leq \psi(r,0)+2L_\rho\int_0^t\psi(r,s)ds\]
    From here we can use our standard Gr\"onwall's inequality to see that 
    
   \[\psi(r,t)\leq \psi(r,0)e^{2L_\rho t}\]
    Recall that $\psi(r,0)\leq rL_0$ and so 
    
   \[|u(x,t)-u(x+h,t)|\leq |h|L_0e^{2L_\rho t}.\] 
    Therefore, for each $t$, $u(\cdot, t)$ is Lipschitz continuous for any finite time and the global Lipschitz constant is $L_0e^{2L_\rho T}$.
    \end{proof}

    In this section we have shown that this model has unique solutions and, for particular $\rho$, those solutions exist globally and are as regular as we can expect. Unlike local diffusion models, nonlocal diffusion problems do not exhibit a smoothening of initial data so we suspect that to gain higher regularity, the regularity of the initial data would have to be increased.   
    
  \section{Analytical Examples}\label{SEC:AnalyticalExamples}
    Having shown that solutions to the IVP without boundary conditions exist and are unique, we now turn our attention to the behavior and properties of solutions. We begin by showing several examples wherein we can write down solutions easily.
    
    \begin{example}[Unstructured Coordination]\label{EX:UnstructuredCoord}
    Consider a bounded domain $\Omega$. If $\rho=-\frac{z^2}{2}$ and $K(x,y)=\frac{1}{Vol(\Omega)}$, we call this the continuous version of the unstructured coordination game (because every player interacts with every other player equally). In this case, the solution can be written down for any $u_0\in C^0_b(\Omega)$.
    
   \[u(x,t)=e^{-t}\left(u_0(x)-\fint_\Omega u_0(y)dy\right)+\fint_\Omega u_0(y)dy\] 
    
    \end{example}
    \begin{proof}
        Observe that our nonlocallity $g$ reduces to 
        \begin{equation*}
            g[u]=\frac{1}{Vol(\Omega)}\int_\Omega u(y,t)-u(x,t)dy=\fint_\Omega u(y,t)dy-u(x,t)
        \end{equation*}
        Once we show that $\fint_\Omega u(y,t)dy$ is constant in time we can solve point-wise as an ODE. Observe that 
        \begin{equation*}
        \begin{split} 
            u(x,t)&=u_0(x)+\int_0^t \left[\fint_\Omega u(y,s)dy-u(x,s)\right]ds\\
            \fint_\Omega u(x,t)dx&=\fint_\Omega u_0(x)dx+\fint_\Omega\int_0^t\left[\fint_\Omega u(y,s)dy-u(x,s)\right]dsdx\\
           \fint_\Omega u(x,t)dx&=\fint_\Omega u_0(x)dx +\int_0^t \left[\fint_\Omega \fint_\Omega u(y,s)dydx-\fint_\Omega u(x,s)dx\right]ds
            \end{split}
        \end{equation*}
        Note that $\fint_\Omega \fint_\Omega u(y,s)dydx=\fint_\Omega u(y,s)dy\fint_\Omega dx=\fint_\Omega u(y,s)dy$ and observe that this implies $\fint_\Omega u(x,t)dx=\fint_\Omega u_0(x)dx$ for all $t$. Therefore we can write our nonlocal equation as
        
       \[u_t(x,t)+u(x,t)=\fint_\Omega u_0(y)dy\] 
        Now it is a simple exercise in ODEs to see that the solution to this is 
        
       \[u(x,t)=e^{-t}\left( u_0(x)-\fint_\Omega u_0(y)dy\right)+\fint_\Omega u_0(y)dy\]

        In particular, if we adjust the initial strategy profile so that it has $\fint_\Omega u_0(y)dy=0$ then $u(x,t)=e^{-t}u_0(t)$. 
    \end{proof}
    The example is illustrative because it validates our model with the expected result in the discrete case. It is easy to see that in the unstructured case here the only equilibrium is a consensus equilibrium. Later in section \ref{SEC:ModelingResults} we will be able to extend result in the unstructured case and say that the only stationary solution (and thus the only equilibrium strategy profile) is the consensus solution whenever $\rho'$ is nonzero away from zero. This is consistent with the result of \cite{Kandori1993}, which found that in the case where all individuals interact, the consensus equilibrium is the only stable equilibrium. The consistency breaks down in the case that $\rho'$ is compactly supported, but in the discrete case all strategies are comparable and so the compact support of $\rho'$ does not have a discrete analog.

    \begin{example}[Unstructured Anti-coordination]\label{EX:UnstructuredAnticoord}
        Consider a  bounded domain $\Omega$ with $\rho(x)=\frac{z^2}{2}$ and $K(x,y)=\frac{1}{Vol(\Omega)}$. The initial value problem with initial date $u_0\in C^0_b$ will have the solution
        
       \[u(x,t)=e^{t}\left( u_0(x)-\fint_\Omega u_0(y)dy\right)+\fint_\Omega u_0(y)dy\]
        
    \end{example}
    Example \ref{EX:UnstructuredAnticoord} is an anticoordination game in the sense that it is the opposite of the coordination game in example \ref{EX:UnstructuredCoord}. It is equivalent to considering the solutions to the coordination game in backward time and for this reason, the solution is immediate. Observe that in this example the coordination condition \eqref{coordinationcondition} is not met and thus the solution does not abide by a maximum principle. However, the solution does exist globally in time.

    \begin{example}[Structured asymmetric-Toeplitz game]\label{EX:Asymmetric}
    Consider a domain $\Omega$ with a Kernel $K\in C^0_b(\Omega;L^1(\Omega))$. If $\rho(z)=cz$ (and thus the game is a dis-coordination game) then the solution is
    
   \[u(x,t)=u_0(x)+tc\|K(x,\cdot)\|_{L^1(\Omega)}.,\]
    \end{example}
    \begin{proof}
        It is easy to see that the nonlocality in this case becomes
        
       \[g[u]=\int_\Omega cK(x,y)dy\] so finding the solution to the IVP is trivial. 
    \end{proof}

    In example \ref{EX:Asymmetric}, we see a glimpse of different behaviors present in Toeplitz games, In example \ref{EX:UnstructuredCoord}, solutions remain bounded and always converge to the consensus equilibrium. Example \ref{EX:UnstructuredAnticoord} is simply the backward time solution of the coordination game and so any constant solution represents an unstable equilibrium but starting from any non-constant initial data will result in a solution growing without bound. In example \ref{EX:Asymmetric} we have a game wherein every player wants to be as far above (or below) the population average as possible so the strategy profile will increase (or decrease) monotonically depending on the kernel $K$. If we combine these elements we can get a situation wherein players are acting under coordination but also seeking to be above-average. 

    \begin{example}[unstructured coordination and advection]\label{EX:UnstructuredAsymmetricCoord}
        Let $\Omega$ be a compact domain. If $K(x,y)=\frac{1}{Vol{\Omega}}$ and $\rho=\frac{-z^2}{2}+cz$ then the solution will be
        
       \[u(x,t)=e^{-t}(u_0(x)-\bar{u_0})+\bar{u_0}+ct\] where $\bar{u_0}=\fint_\Omega u_0(y)dy$
    \end{example}
    \begin{proof}
        Again the nonlocality will be
        
       \[g[u]=\frac{1}{Vol\Omega}\int_\Omega u(y,t)-u(x,t)+cdy\]
       Let $w$ be a general from example \ref{EX:UnstructuredCoord} with $K$ as described and let $v=ct$. $v$ is only a vertical translation so $w(x)-w(y)=(w(x)+v(x))-(w(y)+v(y))$. Let $u=w+v$ and note that $w(x,t)-w(y,t)=u(x,t)-u(y,t)$. Thus 

       \begin{equation*}
       \begin{split}u_t&=w_t+v_t\\
       &=\frac{1}{Vol \Omega}\int_\Omega w(y,t)-w(x,t)dy+\frac{1}{Vol \Omega}\int_\Omega c dy \\
       &=\frac{1}{Vol\Omega}\int_{\Omega}u(y,t)-u(x,t)+c dy\\
       &= g[u]
       \end{split}
       \end{equation*}
        Therefore we can easily see that the solution to the IVP is
        
       \[u(x,t)=e^{-t}\left( u_0(x)- \bar{u_0}\right)+ \bar{u_0}+ct.\]
    \end{proof}

    It appears as though we have combined solutions from examples \ref{EX:UnstructuredCoord} and \ref{EX:Asymmetric}. However, this linear combination of solutions only works because solutions are still solutions under vertical translation. This is the same thing as saying if $u_t=g[u]$ and $w=u+ct$ then $w_t=g[u]+c$. In general, linear combinations of solutions do not solve linear combinations of IVPs. 
    
\section{Numerical Results}\label{SEC:NumericalResults}
    The examples in the previous subsections help us to understand the kinds of behavior we may expect from these kinds of games but we cannot, in general, solve the IVP analytically. However, because of theorem \ref{THM:Regularity}, we can use numerical methods to find solutions, at least in the coordination case. 

    Let us begin on the unit square $\Omega=\prod_{i=1}^n[0,1]$ in $\mathbb{R}^n$ and let $\Omega_T=\Omega\times(0,T)$. If $u_t=g[u]$ on $\Omega_T$ and $u(x,0)=u_0\in C^{0,1}(\Omega)$ then we will try to approximate $u$ with the grid function $w\in \mathcal{V}(\Omega_T^{(h,\tau)})$. Here $\mathcal{V}$ is the set of grid functions which are defined on a discretization of $\overline{\Omega_T}$, 
    
    \begin{equation}\label{hypercube}
    \Omega_T^{(h,\tau)}=\Omega^{h}\times \{\tau l\}_{l=0}^\frac{T}{\tau}=\prod_{i=1}^n\left(\{hk\}_{k=0}^\frac{1}{h}\right)_i\times \{\tau l\}_{l=0}^{T/\tau}.
    \end{equation}
    In particular the set of grid functions we are interested in are $\mathcal{V}(\Omega^{(h,t)}_T)=\{v; v:\Omega^{(h,t)}_T\rightarrow \mathbb{R}\}$ and $\mathcal{V}(\Omega^h)=\{v;v:\Omega^h\to \mathbb{R}\}$ These grid function spaces are different ways of imagining $\mathbb{R}^{\frac{1}{h}^n\cdot \frac{T}{\tau}}$ and $\mathbb{R}^{\frac{1}{h}^n}$ respectively. We use this reimagining so that the comparisons between $v\in \mathcal{V}(\Omega^{(h,\tau)}_T)$ and $u\in C^{0,1}(\Omega_T)$ are more natural. 
    
    Let $\pi^h:C^{0}(\Omega)\rightarrow \mathcal{V}(\Omega^h)$ be the operator which takes a function on $\Omega$ and returns the grid function which is equal to the input function at every point of the grid $\Omega^h$. In the present study we will deal only in the case that $K\in C^0_b(\Omega;C^{0,1}_0)$ so $\pi^hK(x,\cdot)\in \mathcal{V}(\Omega^h)$ for every $x\in \Omega^h$. A generalization is possible but is not immediately necessary for the main results of the paper. There are some cases of $K\in C^0_b(\Omega;L^1(\Omega))$ for which this method is not appropriate. (For example $K(x,y)=d(x,\partial \Omega)^{-n}J(\frac{x-y}{d(x,\partial\Omega)})$ where $J\in L^1(B(0,1))$).

    
    Let $w(\cdot, 0)=\pi^h u_0$ on $\Omega^h$, and compute for an $\mathbf{x}\in \Omega^h$,
    \begin{equation}\label{ForwardEuler}
    w(\mathbf{x},t_{i+1})=w(\mathbf{x},t_i)+\tau \sum_{\mathbf{y}\in ^-\Omega^h}K(\mathbf{x},\mathbf{y})\rho'(w(\mathbf{x},t_i)-w(\mathbf{y},t_i))h^n
    \end{equation}
    where $^-\Omega^h=\prod_{i=1}^n(\{hk\}_{k=0}^{\frac{1}{h}-1})_i$. We will show for this particular domain that the method \eqref{ForwardEuler} is consistent and convergent.
  
    
    \begin{lemma}[Consistency of Forward Euler Method]\label{LEM:Consistency}
        The numerical scheme \eqref{ForwardEuler} is consistent to order $\tau+h$ for the nonlocality \eqref{nonlocality} with a bounded $w\in C^{1,1}$ with no boundary data when $K\in C^0_b(\Omega;C^{0,1}_0)$ with a uniform bounds on $\|K(x,\cdot)\|_\infty$ and on the Lipschitz constant for $K(x,\cdot)$.
    \end{lemma}
    \begin{proof}
          First we compute the error of the right-hand numerical quadrature
          
        \begin{equation}\label{quadrature}\mathcal{G}^h[w](\mathbf{x},t_i):=\sum_{\mathbf{y}\in ^-\Omega^h}K(\mathbf{x},\mathbf{y})\rho'(w(\mathbf{x},t_i)-w(\mathbf{y},t_i))h^n\end{equation}
        for a bounded $w\in C^{0,1}$. Let $\mathbf{y}\in \,^-\Omega^h$ and observe that in the hyperrectangle $\omega_{\mathbf{y}}:=\prod_{k=1}^n[\mathbf{y}_k,\mathbf{y}_k+h]$, $|w(y,t_i)-w(\mathbf{y},t_i)|\leq Lc_nh$ for all $y\in \omega_{\mathbf{y}}$, where $L$ is the Lipschitz constant for $w$ and $c_n$ is a constant depending on the dimension of $\Omega$.  

        As such, for any $y\in \omega_{\mathbf{y}},$ $\mathbf{x}\in \Omega^{h},$ and $\mathbf{y}\in \,^-\Omega^h$, we have $|\rho'(w(\mathbf{x},t_i)-w(y,t_i))-\rho'(w(\mathbf{x},t_i)-w(\mathbf{y},t_i))|\leq L_\rho Lc_nh$ where $L_\rho$ is the Lipschitz constant for $\rho'$ on the interval containing the compact range of $w$, $[-2\|w\|_\infty,2\|w\|_\infty]$. By assumption,  $\|K(\mathbf{x},\cdot)\|_\infty\leq B<\infty$ uniformly over $\Omega^h$ and has a Lipschitz constant, $L_K<\infty$ which is an appropriate Lipschitz constant for $K(x,\cdot) \, \forall x\in \Omega$.  This, along with the fact that $\rho'$ attains its maximum on $[-2\|u\|_\infty,2\|u\|_\infty]$, which we call $W$,  allows us to do the following computation. First note that $K(\mathbf{x},\mathbf{y})\rho'(u(\mathbf{x},t_i)-u(\mathbf{y},t_i))h^n=\int_{\omega_y}K(\mathbf{x},\mathbf{y})\rho'(u(\mathbf{x},t_i)-u(\mathbf{y},t_i))dy$. Thus we write,
        \begin{equation*}
            \begin{split}
                \bigg|\int_{\omega_\mathbf{y}}&K(\mathbf{x},y)\rho'(u(\mathbf{x},t_i)-u(y,t_i))dy-K(\mathbf{x},\mathbf{y})\rho'(u(\mathbf{x},t_i)-u(\mathbf{y},t_i))h^n\bigg|\\
                &\leq\int_{\omega_\mathbf{y}}\bigg|K(\mathbf{x},y)\rho'(u(\mathbf{x},t_i)-u(y,t_i))dy-K(\mathbf{x},\mathbf{y})\rho'(u(\mathbf{x},t_i)-u(\mathbf{y},t_i)\bigg|dy\\
                &\leq \int_{\omega_\mathbf{y}}|K(\mathbf{x},y)-K(\mathbf{x},\mathbf{y})||\rho'(u(\mathbf{x},t_i)-u(y,t_i))|dy\\
                &\quad \quad +\int_{\omega_{\mathbf{y}}}|K(\mathbf{x},\mathbf{y})||\rho'(u(\mathbf{x},t_i)-u(y,t_i))-\rho'(u(\mathbf{x},t_i)-u(\mathbf{y},t_i))|dy\\
                &=:I_1+I_2
            \end{split}
        \end{equation*}
        As we have before, we will deal with each integral separately. The second integral is controlled by $I_2\leq L_\rho Lc_nh\int_{\omega_\mathbf{y}}|K(\mathbf{x},\mathbf{y})|dy\leq B L_\rho Lc_nh^{n+1}$. Now notice that for $I_1$
          \begin{equation*}
              \begin{split}
                  I_1&\leq W\int_{\omega_\mathbf{y}}|K(\mathbf{x},y)-K(\mathbf{x},\mathbf{y})|dy\\
                    &\leq W\int_{\omega_\mathbf{y}}L_Khdy\\
                    &\leq WL_Kh^{n+1}
                \end{split}
          \end{equation*}
          Thus the original difference is controlled by 
          
         \[C_1h^{n+1}:=(WL_K+BL_\rho Lc_n)h^{n+1}\]
        Now when we sum across every $\mathbf{y}\in ^-\Omega^h_T$ we see that 
          \begin{equation*}
              |g[w](\mathbf{x},t_i)-\mathcal{G}^h[w](\mathbf{x},t_i)|\leq C_1h^{n+1}\frac{1}{h^{n}}
          \end{equation*}
        Therefore we have shown that 
        
       \[g[w]=\mathcal{G}^h[w](\mathbf{x},t_i)+\mathcal{O}(h)\]

        Now, suppose $w$ is continuously differentiable in time, it is easy to see that the forward difference 
        
       \[\frac{\partial}{\partial t}w(\mathbf{x},t_i)=\frac{w(\mathbf{x},t_i)-w(\mathbf{x},t_{i+1})}{\tau}+\mathcal{O}(\tau).\] Thus we can say that the numerical scheme \eqref{ForwardEuler} is consistent to order $\tau+h$ for functions $w$ which are bounded, $C^1$ in time, and $C^{0,1}$ in space. 
    \end{proof}
    For a linear problem, the consistency from lemma \ref{LEM:Consistency} and stability from a discrete maximum principle would complete the proof of convergence. The discrete maximum principle holds for this method and is proved in the appendix \ref{APP:AdditionalNumericalResults}, however, it only provides a sanity check that this method makes sense when $\tau$ and $h$ are sufficiently small. Because the equation is nonlinear, we cannot use the Lax principle and lemma \ref{LEM:DiscreteMaximumPrinciple} to prove convergence. Instead, we will recapitulate in the discrete case the argument about the Lipschitz continuity of $\mathcal{G}^h$ in $\Omega^{(h,\tau)}_T$. 

    \begin{lemma}[Lipshitz continuity of $\mathcal{G}^h$]\label{LEM:DiscreteLipschitz}
        For any  bounded (in the sup norm sense) subset of $X_R:=\{w\in \mathcal{V}(\Omega^{(h,\tau)}_T);\|w\|_{\ell^\infty}\leq R\}$, there is a $C^\mathcal{G}\geq 0$ such that for every $w_1, w_2\in X_{R}$
        
       \[\|\mathcal{G}[w_1](\cdot, t_i)-\mathcal{G}[w_2](\cdot, t_i)\|_{\ell^\infty\left(\Omega^h\right)}\leq C^\mathcal{G}\|w_1(\cdot, t_i)-w_2(\cdot, t_i)\|_{\ell^\infty\left(\Omega^h,\right)}\]
        so long as $K\in C^0_b(\Omega;C^{0,1}_0)$ with a uniformly bounded $\|K(x,\cdot)\|_\infty$ and $\mathcal{G}$ is defined as in \eqref{quadrature}.
    \end{lemma}

    \begin{proof}
        Consider $w_1, w_2\in X_R$ and observe that $\|\mathcal{G}[w_1]-\mathcal{G}[w_2]\|_{\ell^\infty}$
        \begin{equation*}
        \begin{split}
            &\leq \max_{\mathbf{x}\in \Omega^h}\left\{\sum_{\mathbf{y}\in ^-\Omega^h}|K(\mathbf{x},\mathbf{y})||\rho'(w_1(\mathbf{x},t_i)-w_1(\mathbf{y},t_i))-\rho'(w_2(\mathbf{x},t_i)-w_2(\mathbf{y},t_i))|h^n\right\}\\
            &\leq \max_{\mathbf{x}\in \Omega^h}\left\{\|\rho'(w_1(\mathbf{x},t_i)-w_1(\cdot,t_i))-\rho'(w_2(\mathbf{x},t_i)-w_2(\cdot,t_i))\|_{\ell^\infty(\Omega^h)}\sum_{y\in^-\Omega^h}|K(\mathbf{x}-\mathbf{y})|h^n\right\}
            \end{split}
        \end{equation*}
        Notice that, because the domain in question has volume $vol(\Omega)=1$, we have \begin{equation*}
        \begin{split}
        \sum_{\mathbf{y}\in^-\Omega^h}|K(\mathbf{x},\mathbf{y})|h^n &\leq \sum_{\mathbf{y}\in \Omega^{h}}\|\pi^hK(\mathbf{x},\mathbf{y})\|_{\ell^\infty(\Omega^{h})}h^n\\
        &\leq\|\pi^h K(\mathbf{x},\cdot)\|_{\ell^\infty(\Omega^h)}\sum_{\mathbf{y}\in^-\Omega^h}h^n\\
        &\leq \|K(\mathbf{x},\cdot)\|_{\ell^\infty(\Omega^h)}vol(\Omega)\\
        &\leq B:=\sup_{\mathbf{x}\in\Omega^h}\|K(\mathbf{x},\cdot)\|_{\ell^\infty(\Omega^h)}
        \end{split}
        \end{equation*}
        for all $\mathbf{x}\in \Omega^h$ and so we proceed with a very similar argument about $\rho'$ as in lemma \ref{LEM:Lipschitzg}.

        \begin{equation*}
            \begin{split}
                \|\rho'(w_1(\mathbf{x},t_i)-&w_1(\cdot, t_i))-\rho'(w_2(\mathbf{x},t_i)-w_2(\cdot,t_i))\|_{\ell^\infty(\Omega^h)}\\
                &\leq L_\rho\|w_1(\mathbf{x},t_i)-w_1(\cdot,t_i)-w_2(\mathbf{x},t_i)+w_2(\cdot,t_i)\|_{\ell^\infty(\Omega^h)}\\
                &\leq L_\rho\left(|w_1(\mathbf{x},t_i)-w_2(\mathbf{x},t_i)|+\|w_1(\cdot,t_i)-w_2(\cdot,t_i)\|_{\ell^\infty(\Omega^h)}\right)
            \end{split}
        \end{equation*}
        So naturally we see, as before, that
        \begin{equation*}
        \begin{split}
        \max_{\mathbf{x}\in \Omega^h}\|\rho'(w_1(\mathbf{x},t_i)-&w_1(\cdot, t_i))-\rho'(w_2(\mathbf{x},t_i)-w_2(\cdot,t_i))\|_{\ell^\infty(\Omega^h)}\\
        &\leq 2L_\rho\|w_1(\cdot, t_i)-w_2(\cdot, t_i)\|_{\ell^\infty(\Omega^h)}
        \end{split}
        \end{equation*}
        Therefore, we get the result that
        
       \[\|\mathcal{G}[w_1](\cdot ,t_i)-\mathcal{G}[w_2](\cdot, t_i)\|_{\ell^\infty(\Omega^h)}\leq 2L_\rho M\|w_1(\cdot, t_i)-w_2(\cdot, t_i)\|_{\ell^\infty(\Omega^h)}\]
        Let $C^\mathcal{G}=2L_\rho M$ to complete the proof. Again $L_\rho$ may depend on $R$ so $C^\mathcal{G}$ will depend on $R$. 
    \end{proof}
    With these two results, we can show that the numerical scheme \eqref{ForwardEuler} in convergent in the case that $K\in C^0_b(\Omega;C^{0,1}_0)$ with uniformly bounded $\|K(x,\cdot)\|_\infty$ and with a global Lipschitz constant. The general case wherein $K\in C^0_b(\Omega;L^1(\Omega))$ requires a choice of approximate $K$ which are bounded and defined everywhere to manage the case in which $K(x,\cdot)$ has a singularity for some $x\in \Omega$ which makes this method inappropriate in some cases. There are more sophisticated numerical methods for non-local equations which are equipped to handle the more general case \cite{D’Elie2020} but we do not use them in the analysis in section \ref{SEC:ModelingResults}. For this reason we only consider the finite difference method in this nonlinear setting, which mildly extends the results from \cite{Du2019} for the linear setting.
    
    \begin{theorem}[Convergence of the Forward Euler Scheme]\label{THM:Convergence}
    In a domain with discretization \eqref{hypercube}, the IVP $u_t=g[u]$ on $\Omega$, with $u(x,0)=u_0\in C^{0,1}_b(\Omega)$, and with $\rho$ satisfying \eqref{coordinationcondition}, the numerical scheme \eqref{ForwardEuler} is convergent in the case where $K\in C^0_b(\Omega;C^{0,1}_0)$ with uniformly bounded $\|K(x,\cdot)\|_\infty$ and a uniform global Lipschitz constant.
    \end{theorem}
        
    \begin{proof}
        Let $\pi_{h,\tau}:C^0(\Omega_T)\rightarrow \mathcal{V}(\Omega^{(h,\tau)}_T)$ discretize $u$, a solution of the IVP $u_t=g[u]$ with $u(x,0)=u_0\in C^{0,1}_b(\Omega)$. Let $w$ be a grid function which solves the numerical scheme \eqref{ForwardEuler} with $w(\mathbf{x},0)=\pi_h u_0$. Further let $e=\pi_{h,\tau} u- w$. Now notice that 
        \begin{equation}
            \begin{split}
            \frac{e(\mathbf{x},t_{i+1})-e(\mathbf{x},t_{i})}{\tau}&=\frac{u(\mathbf{x},t_{i+1})-u(\mathbf{x},t_{i})}{\tau}- \mathcal{G}^h[w](\mathbf{x},t_i)\\
            &=u_t(\mathbf{x},t_i)-\mathcal{G}^h[w](\mathbf{x},t_i)+\mathcal{O}(\tau)\\
            &=g[u](\mathbf{x},t_i)-\mathcal{G}^h[w](\mathbf{x},t_i)+\mathcal{O}(\tau)\\
            &=\mathcal{G}^h[u](\mathbf{x},t_i)-\mathcal{G}^h[w](\mathbf{x},t_i)+\mathcal{O}(\tau+h)
            \end{split}
        \end{equation}
        by lemma \ref{LEM:Consistency}. $\mathcal{G}^h$ is not linear so the stability of the method does not complete the proof. Instead, we use lemma \ref{LEM:DiscreteLipschitz} to show that
        \begin{equation*}
            \begin{split}
        e(\mathbf{x},t_{i+1}) &=e(\mathbf{x},t_{i})+\tau\left(\mathcal{G}^h[u](\mathbf{x},t_i)-\mathcal{G}^h[w](x,t_i)+\mathcal{O}(\tau+h)\right)\\
        \|e(\cdot, t_{i+1})\|_{\ell^\infty(\Omega^h)}&\leq\|e(\cdot, t_{i})\|_{\ell^\infty(\Omega^h)}+ \tau \|\mathcal{G}^h[u](\mathbf{x},t_i)-\mathcal{G}^h[w](x,t_i)\|_{\ell^\infty(\Omega^h)}+\tau|\mathcal{O}(\tau+h)|\\
        &\leq\|e(\cdot, t_{i})\|_{\ell^\infty(\Omega^h)}+ \tau C^\mathcal{G}\|\pi^h u(\mathbf{x},t_i)-w(x,t_i)\|_{\ell^\infty(\Omega^h)}+\tau|\mathcal{O}(\tau+h)|\\
        &\leq (1+\tau C^\mathcal{G})\|e(\cdot, t_i)\|_{\ell^\infty(\Omega^h)}+\tau|\mathcal{O}(\tau +h)|
        \end{split}
        \end{equation*}

        We now employ Gr\"onwall's lemma in the discrete forward difference case to say that
        \begin{equation*}
            \|e(\cdot, t_{i+1})\|_{\ell^\infty(\Omega^h)}\leq (1+\tau C^\mathcal{G})^i\|e(\cdot, 0)\|_{\ell^\infty(\Omega^h)}+\frac{1}{C^\mathcal{G}}((1+\tau C^\mathcal{G})^i-1)C(\tau+h)
        \end{equation*}
        Naturally, $e(\cdot, 0)\equiv 0$ and for finite time $T$ we have a maximum $i\leq T/\tau$ so we can say that
        \begin{equation*}
            \|e\|_{\ell^\infty(\Omega^{(\tau,h)}_T)}\leq \tilde{C}((1+\tau C^\mathcal{G})^\frac{T}{\tau}-1)(\tau+h)
        \end{equation*}
        Observe that $C^\mathcal{G}$ does not depend on $h$ or $\tau$. It only depends on $R$ but because of lemma \ref{LEM:MaximumPrinciple} and its discrete analogue (lemma \ref{LEM:DiscreteMaximumPrinciple}), $R$ is the same at each time step and thus so is $C^\mathcal{G}$. Also observe that $\lim_{\tau\rightarrow 0}(1+\tau C^\mathcal{G})^\frac{T}{\tau}$ exists and so this quantity is bounded for any sufficiently small $\tau$ and indeed as $(\tau, h)\rightarrow \mathbf{0}$, $\|e\|_{\ell^\infty}\rightarrow 0$.
    \end{proof}

    This proof was done for the unit cube but each part is easily generalizable to any bounded domain with constants which may depend on the volume of the domain itself.  Now that we are certain that this numerical scheme appropriately approximates solutions to the IVP, we can show several examples.
    \begin{figure}
        \centering
        \includegraphics[width=0.5\linewidth]{ICfigure.png}
        \includegraphics[width = 0.49\linewidth]{NCPTfigure.png}
        \caption{The Initial condition (\textbf{left}) and solution after $T=10$ (\textbf{right}) to the IVP $u_t=g[u]$ approximated by the numerical scheme \eqref{ForwardEuler}. Here the recognition function has non-zero derivative on all of $\mathbb{R}\setminus\{0\}$}
        \label{fig:NCPT}
    \end{figure}
    \begin{figure}
        \centering
        \includegraphics[width=0.5\linewidth]{ICfigure.png}
        \includegraphics[width = 0.49\linewidth]{CPTfigure.png}
        \caption{The Initial condition (\textbf{left}) and solution after $T=10$ (\textbf{right}) to the IVP $u_t=g[u]$ approximated by the numerical scheme \eqref{ForwardEuler}. Here the recognition functions is identically 0 outside $B_{1/4}(0)$.}
        \label{fig:CPT}
    \end{figure}

    Interrogating the system with numerical methods allow us to observe some curious properties of solutions to the coordination equation. One such interesting observation is the qualitatively different behavior seen for solutions which depends on the support of $\rho$ as in figures \ref{fig:NCPT} and \ref{fig:CPT}. This will be discussed further in the modeling results. 

 \section{Modeling Results}\label{SEC:ModelingResults}
    Having shown that the model is well posed, solutions exist, and that solutions can be approximated through simple numerical methods, we turn our attention now to what this model may reveal about coordination in continuous settings. The first, and most striking observation is the apparent discontinuities which emerge in the limit at $t\to \infty$ when $\rho'$ has compact support (as in Fig \ref{fig:CPT}). It has not been proven that solutions to the IVP will converge, even pointwise to a limit, but we do know that, if they do converge, they will clearly converge to a solution to the problem $g[u]=0$ in $\Omega$. Without imposing boundary data, solutions to this problem obviously exist (e.g. $u\equiv 0$). The existence of non-trivial solutions and solutions with boundary data are not discussed in the present study. We will, however, discuss several results about stationary solutions and present some results from numerical experiments. 

    \subsection{Stationary Solutions}\label{SUBSEC:StationarySolutions}
    Recall that, as the game was introduced in section \ref{SEC:Extension}, we are not only interested in the dynamic results but in fact may wish to consider the classical game with no time component. Results about stationary solutions in the dynamic game can, unsurprisingly, reveal more effective ways of searching for Nash equilibria in the classical game. If $u$ is a Nash equilibrium in the classical game, by definition it will have the property that 
    \begin{equation}\label{NECondition}
    J[u]:=\inf_{(x,s)\in \Omega\times \mathbb{R}}\{w(x|u)-w(x|u+t\chi_{\{x\}})\}=0.\end{equation}
    This is exactly the condition that each player is playing their best response. 
    \begin{proposition}\label{PROP:StationarySolutions}
        If $u$ is a Nash Equilibrium to the game with players in $\Omega\subset\mathbb{R}^n$, strategies in $\mathbb{R}$ and payoffs as in \eqref{ctsStratPayoff}, then $u$ is necessarily a stationary solution in the system $u_t=g[u]$.
    \end{proposition}
    \begin{proof}
        In the same way as in the proof of proposition \ref{PROP:ModelWellFounded} let $S_x(h):=\int_\Omega K(x,y)\rho(u(x)+h-u(y))dy$ and note that, so long as $\rho\in C^{1,1}(\mathbb{R})$ then $S_x(h)\in C^{1,1}(\mathbb{R})$ and $\frac{d}{dh}S_x(0)=\int_\Omega K(x,y)\rho'(u(x,t)-u(y,t))$. Because $u(x)$ is a best response to $u$, it is certain that $S_x(h)$ attains a global maximum at $h=0$ and because the strategic domain is unbounded, we know that $\frac{d}{dh}S_x(0)=\int_\Omega K(x,y)\rho'(u(x)-u(y))dy=0$. This is true for every $x$ and thus $g[u]=0$ in $\Omega$.
    \end{proof}

    It should be noted that the opposite direction does not hold. It is easy to construct a stationary solution which is not a Nash equilibrium. 

    
    \begin{example}
        Consider a coordination game where $\rho\in C^{1,1}(\mathbb{R})$ satisfies $\rho(0)>0, \rho'(0)=0,$ and $supp(\rho')\subset[-a,a]$.
        
       \[u=\begin{cases}
        0&x\neq0\\
        2a&x=0
        \end{cases}\]
        is a stationary solution because $g[u]\equiv 0$. However, $u$ is not a Nash equilibrium because $w(0|u)=0$ but when $\tilde{u}=0$ on all of  $\Omega$, (so $\tilde{u}=u-2a\chi_{\{0\}}$) then $w(0|\tilde{u})=\rho(0)\|K(0,\cdot)\|_{L^1(\Omega)}$. Thus we have that $J[u]\leq w(x|u)-w(x|\tilde u)= -\rho(0)\|K(0,\cdot)\|_{L^1(\Omega)}<0$ where $J$ is defined as in \eqref{NECondition}.
    \end{example}

    This means that understanding stationary solutions will inform our understanding of the classical game, even if we cannot connect,  rigorously, our understanding of the stationary solutions to the dynamics of the IVP we have been studying. In the following theorem we characterize stationary solutions in the case that $\Omega$ is bounded and there is no boundary data whenever $K$ is supported on all of $\Omega$.

     \begin{theorem}[Stationary Solutions when $supp(\rho')$ is compact]\label{THM:StationarySolution}
        Let $\Omega$ be a bounded domain in $\mathbb{R}^n$. Let $K(x,y)\in C_b^0(\Omega;L^1(\Omega))$ so that for any $x$, $\Omega\subseteq supp(K(x,\cdot))$ and $\lambda \leq K(x,y)$. Finally, let $\rho'$ satisfy \eqref{coordinationcondition} and have support $(-a,a)\subset \mathbb{R}$, with the assumption that $\rho'$ has only one zero in this interval at $z=0$, see figure \ref{fig:rhoprime}. If $u$ is a solution to $g[u]=0$ in $\Omega$ and $u$ is bounded, then the image, $u(\Omega)$, is a finite set of points separated by at least $|a|$ except possibly at a set of measure 0.
    \end{theorem}

    \begin{figure}
        \centering
        \begin{tikzpicture}
            \draw[<->](-7,0)--(7,-0);
            \draw[<->](0,-3)--(0,3);

            \node (a) at (-0.5,3){$\rho'(z)$};
            \node (b) at (7,-0.5){$z$};

            \draw(-5,0)--(-5,-0.25);
            \draw(5,0)--(5,0.25);
                
            \node (c) at (-5,-0.5){$-a$};
            \node (f) at (5,0.5){$a$};
        

            \draw plot [smooth] coordinates {(-5,0) (-4,0.5) (-3,2) (-2,1) (-1,1.5) (0, 0) (1,-0.75) (1.5,-0.6) (2,-1) (2.5,-1.5) (3,-2) (3.5,-2)(4,-1) (5,0)};
        \end{tikzpicture}
        \caption{A diagram showing an appropriate $\rho'$ for theorem \ref{THM:StationarySolution}. $\rho'$ must be non-zero on $(-a,a)\setminus\{0\}$ and satisfy \eqref{coordinationcondition} so that on $(-a,0)$ $\rho'$ is positive and on $(0,a)$, $\rho'$ is negative. $\rho'$ is Lipschitz continuous but no more regularity is required for this argument.}
        \label{fig:rhoprime}
    \end{figure}

    
    \begin{proof}
        First note that by the assumption on $K$, $C^K\geq \|K(x,\cdot)\|_{L^1(\Omega)}>\lambda Vol(\Omega)=:C_K$ for all $x$. 
        Let $M_0=\sup_{\Omega}u(x)$. Whether or not it is attained, there is a sequence of $x_k$ so that $u(x_k)\to M_0$ and $k\to \infty$. For each $k$ we partition the domain into three parts
        \begin{equation*}
        \begin{split}
            \Omega^+_k&:=\{x\in \Omega, u(x)>u(x_k)\}\\
            \Omega^o_k&:=\{x\in \Omega; u(x)=u(x_k)\}\\
            \Omega^-_k&:=\{x\in \Omega; u(x)<u(x_k)\}
        \end{split}
        \end{equation*}
        and so $g[u]$ is partitioned into three parts
        \begin{equation*}  
            \begin{split} 
                g[u](x_k)& =\underbrace{\int_{\Omega^+_k}K(x_k,y)\rho'(u(x_k)-u(y))dy}_{:=I_k^+\geq 0}+\underbrace{\int_{\Omega^o_k}K(x_k,y)\rho'(u(x_k)-u(y))dy}_{:=I_k^o=0}\\
                &\quad +\underbrace{\int_{\Omega^-_k}K(x_k,y)\rho'(u(x_k)-u(y))dy}_{:=I_k^-\leq0}\end{split}\end{equation*}
        Note that for every $k$, $I^+_k=-I^-_k$ because $g[u]=0$. 

        Let $\beta>0$ and note that we can find a $K_\beta$ so that $u(x_k)>M_0-\frac{\beta}{C^KL_\rho}$ for all $k>K_\beta$ where $L_\rho$ is the Lipschitz constant for $\rho'$. Because $\rho'$ is Lipschitz continuous and $\rho'(0)=0$, if $-\frac{\beta}{C^KL_\rho}<u(x_k)-u(y)<0$ for all $y\in \Omega^+_k$, then $\rho'(u(x_k)-u(y))\leq \frac{L\beta}{C^KL}=\frac{\beta}{C^K}$ for all $y\in \Omega^+_k$. 
        This means that for any $\beta$ we can find a $K_\beta$ so that when $k>K_\beta$,
        \[I_k^+<\|K(x_k,\cdot)\|_{L^1(\Omega)}\frac{L\beta}{C^KL}\leq\frac{C^KL\beta}{C^KL}=\beta. \]
        This means that for every $\beta$ sufficiently small, we must see that there exists a $K$ (indeed the same $K_\beta$ should do) such that $|I_k^-|\leq\beta$ for all $k>K_\beta$.

        Now, consider $\epsilon>0$ sufficiently small and suppose there is an $S\subset\Omega$ with positive measure $\mu(S)>0$ and satisfying the inequality $M_0-a+\epsilon<u(s)<M_0-\epsilon$. Let $K_{\epsilon}$ be the index so that $k>K_{\epsilon}\implies M_0-u(x_k)<\frac{\epsilon}{2}$. Note that when $k>K_{\epsilon}$, $S\subset \Omega ^-_k$. Moreover, because $\frac{\epsilon}{2}\leq u(x_k)-u(s)\leq a-\epsilon$, and $\rho'$ has no zeros in this compact interval, there is a $C_\epsilon$ such that $\rho'(u(x)-u(s))\le C_\epsilon<0$. This means that when $k>K_{\epsilon}$ we can bound $I^-_k$ away from $0$
        \[
        -I^-_k\geq -\int_SK(x,y)\rho'(u(x_k)-u(y))dy\geq -\lambda C_\epsilon \mu(S)>0.
        \]

        This contradicts the fact that $I^-_k\to 0$ as $k\to \infty$ so we conclude that for any subset $S\subset \Omega$ with $M_0-a+\epsilon_0\leq u(S)\leq M_0-\epsilon_0$, we know that $\mu(S)=0$. Restated, we have shown that in the bounded domain $\Omega$, 
        \begin{equation}\label{EQ:StationaryBaseCaseEpsilon} u(s)\geq M_0-\epsilon\quad\text{ or }\quad u(s)\leq M_0-a+\epsilon\quad \text{ almost everywhere.}\end{equation} 

        Importantly, every step taken to arrive at \eqref{EQ:StationaryBaseCaseEpsilon} can still be done when $\epsilon$ is made smaller and so in the limit as $\epsilon\to 0$ we can conclude that
        
        \begin{equation}\label{EQ:StationaryBaseCase} u(s) = M_0\quad\text{ or }\quad u(s)\leq M_0-a\quad \text{ almost everywhere.}\end{equation} 
        To see the convergence argument, observe that if \eqref{EQ:StationaryBaseCase} was not true we would surely have that $|\{s\in \Omega;M_0-a<u(s)<M_0\}|>0$ and we can write this as the countable union 
        \[|\{s\in \Omega;M_0-a<u(s)<M_0\}|=\bigg|\bigcup_{\epsilon=\frac{1}{n}}\{s\in\Omega; \epsilon-a<u(s)<M_0<\epsilon\}\bigg|>0.\]
        In order for this to be positive, it must be positive for at least one $\epsilon$ and this could contradict \eqref{EQ:StationaryBaseCaseEpsilon}.

        
        Now, consider $\{u(x);x\in \Omega, u(x)\neq M_0\}$. If this set has measure 0 then we have completed the proof. If it has positive measure, let $M_1:=\esssup\{u(x);x\in \Omega, u(x)\neq M_0\}$ and note that $M_1\leq M_0-a$. Let $x_k\in \Omega$ so that $u(x_k)\leq M_1$ and $u(x_k)\to M_1$. Now partition the domain into 
        \begin{equation*}
            \begin{split}
            \Omega^{++}_k&:=\{x\in \Omega; u(x)\geq u(x_k)+a\}\\
                \Omega^+_k&:=\{x\in \Omega, u(x_k)<u(x)<u(x_k)+a\}\\
            \Omega^o_k&:=\{x\in \Omega; u(x)=u(x_k)\}\\
            \Omega^-_k&:=\{x\in \Omega; u(x)<u(x_k)\}
            \end{split}
        \end{equation*}
        and again we can split the integral into
        \begin{equation*}
            \begin{split}
                g[u](x_k)& =\underbrace{\int_{\Omega^{++}_k}K(x_k,y)\rho'(u(x_k)-u(y))dy}_{:=I_k^{++}= 0}+\underbrace{\int_{\Omega^+_k}K(x_k,y)\rho'(u(x_k)-u(y))dy}_{:=I_k^+\geq 0}\\
                &\quad+\underbrace{\int_{\Omega^o_k}K(x_k,y)\rho'(u(x_k)-u(y))dy}_{:=I_k^o=0} +\underbrace{\int_{\Omega^-_k}K(x_k,y)\rho'(u(x_k)-u(y))dy}_{:=I_k^-\leq0}
            \end{split}
        \end{equation*}

        In this way we are left with the same situation as before. $I^{++}_k=0$ because when $u(y)\geq u(x_k)+a$ then $u(x_k)-u(y)\leq -a$ and so $\rho'(u(x_k)-u(y))=0$ for all $y\in \Omega_k^{++}$. Thus, we note that $I^+_k=-I^-_k$ for all $k$. Again, for any $\beta$ we can find a $K_\beta$ so that $I^+_k<\beta$ for all $k>K_\beta$. This is because we proved in the first part of the proof that $u(x)\notin (M_0-a,M_0)$ except for possibly at a set of measure $0$. Thus the entire contribution to $I^+_k$ comes from $x\in \Omega^+_k$ with $u(x_k)<u(x)\leq M_1$ and we can repeat our previous argument exactly. Thus $u(x)\notin (M_1,M_1+a)$ except possibly at a set of measure $0$. Because the nonlocality will not see a set of measure $0$, it is an identical argument to extend the result in \eqref{EQ:StationaryBaseCase} to say that in $\Omega$,
        \[u(s)=M_0\quad \text{ or } \quad u(s)=M_1 \quad\text{ or } \quad u(s)\leq M_1-a \quad\text{ almost everywhere}\]

        This argument can be repeated indefinitely but notice that between each of the resulting bands in $u(\Omega)$ there is a gap of at least $a$. This means that, after a finite number of repetitions we will cover the entire range of the bounded solution, $u$. This means that, except for a possible set $Q$ which has measure 0, the image of $u(\Omega\setminus Q)$ is a finite set of points. Further, we can partition $\Omega$ into $A_0,...,A_m$ where $A_k:=\{x\in \Omega; u(x)=M_k\}$ (by the process described above each of these $A_i$ will have positive measure) and represent the solution $u$ as a simple function
        \[u(x)= \sum_{k=0}^mM_k\chi_{A_i=k}(x)\in L^\infty(\Omega)\] 
    \end{proof}

    This result is remarkable and gives us some insight into coordination dynamics at equilibrium. Although we cannot rule out the existence of a set of players with measure zero that do not adhere to one of a finite set of strategies, we can say that, in a weak sense, at equilibrium when no outside influences are acting on the system (i.e. no boundary data or inhomogeneity), the solution will be piecewise constant. If we know apriori that the solution is continuous then we can go even further.
    \begin{corollary}\label{COR:ContinuousStationarySolution}
        If $u$ is a continuous bounded stationary solution on a bounded domain with $\rho$ having only 1 zero in the interior of its support at $z=0$ and $K$ supported on the entire domain,  then $u$ is constant
    \end{corollary}
    
    In order to see continuous and non-constant solutions, at least one of these assumptions must be broken. An easy example is when the boundedness of the solution and of the domain is violated. There are clear examples of solutions to $g[u]=0$ which are unbounded and not piecewise constant. 
    \begin{example}\label{EX:NonconstantStationarySolution}
        Let $k(x,y)=J(|x-y|)$ a radial, translation invariant kernel. Let $\rho$ be even so $\rho'$ is odd. For the domain $\mathbb{R}$, the solution $u(x)=x$ satisfies $g[u]=0$. 
    \end{example}
    \begin{proof}
        $\rho'$ is odd so for any $x$, $\rho'(u(x)-u(y))=\rho'(x-y)$ is odd in $y$. $J$ is even and translation invariant so again for any $x$ it is even in $y$. Thus $g[u](x)=\int_{\Rn}J(x-y)\rho'(x-y)dy=0$ for all $x\in \Omega$. 
    \end{proof}

    Another example may be if $\rho'$ is not supported in an interval around $z=0$. In this case, regardless of $K$ we can construct a continuous non-constant stationary solution
    
    \begin{example}
        $K$ is translation invariant and has support in $B_\delta(0),$ and $\rho$ satisfies 
        \[
            \rho(z) =\begin{cases}
            1&|z|<r_0\\
            0&|z|\geq r_1
        \end{cases}
        \]
          where $r_0 <r_1$ and  $0\leq \rho(z) \leq 1$ when $r_0\leq |z|\leq r_1$.
    than any $u$ which is globally Lipschitz with Lipschitz constant $\frac{r_1}{\delta}$ or which satisfies $\sup(u)-\inf(u)<r_0$ is a stationary solution.
    \end{example}

    These results about stationary solutions are helpful because they give us a way of narrowing down the search for Nash equilibria in the classical game. By proposition \ref{PROP:StationarySolutions} and by Theorem \ref{THM:StationarySolution}, we know that any bounded Nash equilibrium of the game with appropriate $K$ and $\rho'$ on a bounded domain will have a finite image when at most a set of measure 0 is excluded from the domain. We do not expect this result to hold in the case that there are boundary data nor do we expect this to hold if we consider the inhomogeneous problem discussed in section \ref{SEC:Inhomogeneous}.

    From the proof of Theorem \ref{THM:StationarySolution}, we can see that there is upper bound on the number of points present in the image of a stationary solution. Namely, if $b$ is the number of points in the image (which will appear as bands on a bifurcation diagram in the next subsection), $R:=\sup_{x\in\Omega}u(x)-\inf_{x\in\Omega}u(x)$ and $r$ is the radius of $supp(\rho')$, then $b\leq\lfloor\frac{R}{r}\rfloor+1$. Notably this means that if $\rho'$ is supported on all of $\mathbb{R}$ and has only one zero at $z=0$ then a stationary solution will be constant except possibly at a set of measure zero. This upper bound is sharp as we can construct stationary solutions which satisfy $b=\lfloor\frac{R}{r}\rfloor+1$ easily.
    \[u(x)=\sum_{k=1}^{b}rk\chi_{S_k}=\begin{cases}
        r&x\in S_1\\
        2r&x\in S_2\\
        \vdots&\\
        br&x\in S_{b}
    \end{cases}\]
    Where $\{S_i\}_{i=1}^b$ partitions $\Omega$. This result characterizes the stationary solutions but, because of the lack of a time convergence result, we cannot yet unite the dynamics described in diffusion problem $u_t=g[u]$ with their apparent limits rigorously. In order to get some idea of how the time dependent coordination process proceeds towards equilibrium, we use several numerical experiments. 
    \subsection{Numerical Experiments}\label{SUBSEC:NumericalExperiments}
    From the above results about stationary solutions we seek to investigate the behavior of solutions as time tends towards infinity. We conducted several numerical experiments employing the numerical methods described in section \ref{SEC:NumericalResults}. The code used to run the experiments can be found at \cite{McAlister2025}. The first experiment considered the interval $I=[-1/2,1/2]$ and the initial data $u_0(x)=lx$ where $l$ was varied from $0$ to $4$. With a Gaussian kernel $K(x,y)=\frac{1}{s\sqrt{2\pi}}e^\frac{{(x-y)}^2}{-2s^2}$ (where the kernel concentrates around $x$ as $s$ increases) and a compactly supported recognition function $\rho(z) = e^{\frac{-1}{1-(z/r)^2}}$ when $|z|<r$ and $0$ otherwise, the experiment examines the profile of the solution after the IVP has been solved until $t=20$. The system parameters were chosen as $s=0.5$ and $r=0.2$. The results are shown in Figure \ref{fig:NM1}
    \begin{figure}
        \centering
        \includegraphics[width=\linewidth]{NM1TotalHotCrop.png}
        \caption{\textbf{Top} A bifurcation-type diagram showing that as the slope of the initial data increases, the distribution of the solution $u(x,20)$ changes. On the far left, when the slope is low, the image in entirely distributed at $u(x)-=0$. As the initial slope increases there are more bands. The colors on the heat map represent the $log_2$ of the density of $u(x,20)$.
        \textbf{Bottom} Three examples of initial conditions (dotted in blue) and solutions at time $t=20$ (solid in red). The corresponding slices in the bifurcation diagram are labeled in order on the $x$ axis of the heatmap above.}
        \label{fig:NM1}
    \end{figure}

    An interesting observation about the number of bands present at time $t=20$ is the approximate equal spacing between the non-central bands. All of the non-central bands are separated by $0.4$, which is the diameter of the support of $\rho'$. This is unlike the example from subsection \ref{SUBSEC:StationarySolutions} where the minimum distance between bands could be only the radius of the support.  When there are two central bands, they can achieve this minimum separation but will grow apart as $l$ increases. It is exactly when the distance between the bands exceed $0.4$ when the appearance of a single central band between them emerges. Following this pattern we can predict, for a linear initial condition, how many bands there will be. No two bands can be closer than 0.2 (half the support $\rho'$) and no two bands can be further than 0.4. For this setting in particular, at most 2 spaces between bands will be less than 0.4. Therefore we can arrive at the bound that the number of bands $b$ is limited by 
    \[b\leq \frac{R}{2r}+1\]
    where again $R$ is the length of the range of $u_0$ and $r$ is the radius of the support of $\rho'$. This bound is smaller than the upper bound from subsection \ref{SUBSEC:StationarySolutions} which may be related to the stability of these stationary solutions. 

    Another interesting observation which is not explored further in the present study is the fact that the bands seems to start to form from the outside then proceed in. In the case where $l=3.5$ in figure \ref{fig:NM1} we can see that the most central bands have not formed by the time $t=20$. This points to a possible finite propagation speed. As in example \ref{EX:NonconstantStationarySolution}, when the solution $u(x,t)$ is odd about a point $x_0$ and the kernel is even, then $u_t(x_0,t)=0$. In the exact center the solution is odd for all time in each of these experiments. Close to the center the solution ``almost" has odd symmetry and near the edge of the range, the solution has no odd symmetry. For points near the center, the lack of odd symmetry far away is hidden by the fact that the kernel $K$ is small at that distance. For points close to the edge of the range, the immediate lack of odd symmetry results in an immediate collapse into discrete strategies, this collapse ruins the odd symmetry closer to the center and so the collapse propagates to the center (Fig \ref{fig:NM4}).
    \begin{figure}
        \centering
        \includegraphics[width=\linewidth]{NM4}
        \caption{Four time slices of the solution to $u_t=g[u]$ when $u_0=7/2x$ on the domain $\Omega =[-2,2]$. The collapse into discrete strategies starts at the edge of the domain and propagates towards the center.}
        \label{fig:NM4}
    \end{figure}

    We use the next two numerical experiments to test the hypothesis about the minimum band separation. Using a sigmoid function instead of a linear function for the initial data we can investigate the number of bands that emerge while the range of the initial data is constrained. For the same kernel and recognition function as before, we run the same experiment with $u_0=(1+\exp(-lx))$ for $l\in [0,15]$. The results are shown in figure \ref{fig:NM2} 

    \begin{figure}
        \centering
        \includegraphics[width=\linewidth]{NM2TotalHotCrop.png}
        \caption{\textbf{Top} A bifurcation-type diagram showing that as the sigmoid parameter of the initial data increases, the distribution of the solution $u(x,20)$ changes. On the far left, when the sigmoid parameter is low, the image is entirely distributed at $u(x)=0$. As the sigmoid parameter for the initial data increases, there are two then later three bands but as the sigmoid parameter tends towards its maximum, the bands level out. The colors on the heat map represent the $log_2$ of the density of $u(x,20)$.
        \textbf{Bottom} Three examples of initial conditions (dotted in blue) and solutions at time $t=20$ (solid in red). The corresponding slices in the bifurcation diagram are labeled in order on the $x$ axis of the heatmap above.}
        \label{fig:NM2}
    \end{figure}
    Again we notice that the number of bands increases with the range of the initial data. It is straight forward to observe that the relationship from the first numerical experiment holds in this case too. Because the range of the initial data is constrained to $(0,1)$ we never see more than three bands emerge. This is precisely what is predicted from the hypothesis previously stated. 

    The sigmoid example shows us also that it is not a critical gradient threshold that results in discontinuities in the limit. The derivative is greatest in at $x=0$ initially in every example, however, the discontinuities do not always emerge at the center. When the sigmoid parameter is greater than $5$ and there are three bands present, the central band is centered around $x=0$ and the discontinuities are present on either side of this band. Therefore, when trying to predict the location of bands present at equilibrium (assuming a solution converges to an equilibrium) is not as simple as finding where the gradient might surpass a certain threshold. 

    To further test this hypothesis, we vary the support of the recognition function rather than the initial data. In the third numerical experiment, the parameter $r$ in the recognition function 
    \[\rho(z)=\begin{cases}e^\frac{-1}{1-(z/r)^2} & |z|<r\\
    0&|z|\geq r\end{cases}\]
    is varied from $0$ (the limit in which $\rho(z)=\chi_{\{0\}}(z)$) to $0.8$. The results are shown in figure \ref{fig:NM3} 

    \begin{figure}
        \centering
        \includegraphics[width=\linewidth]{NM3TotalHotCrop.png}
        \caption{\textbf{Top} A bifurcation-type diagram showing that as support of the recognition function changes, the distribution of the solution $u(x,20)$ changes. On the far left, when the support of $\rho$ is near $0$, there are many bands. As the support of the recognition function increases, there are fewer and fewer bands until the solution collapses to $u(x)= 0$. The colors on the heat map represent the $log_2$ of the density of $u(x,20)$.
        \textbf{Bottom} Three examples of initial conditions (dotted in blue) and solutions at time $t=20$ (solid in red). The corresponding slices in the bifurcation diagram are labeled in order on the $x$ axis of the heatmap above.}
        \label{fig:NM3}
    \end{figure}

    The results of the experiment are predicted exactly by our hypothesis. Again, in every case, the number of bands is bounded above by $(R/2r)+1$.

    These numerical experiments obviously fall far short of a real proof, and without a time-convergence result such a proof is likely impossible. However, they do force us to consider which stationary solutions are discoverable (i.e. they have a non-trivial basin of stability) and which are not. 

    \section{Towards an Inhomogeneous Problem}\label{SEC:Inhomogeneous}
    Recall that this model is inspired by a game theoretic situation in which two players are interacting and the payoff of their interaction depends only on the ``distance" between their strategies. We can extend this idea to consider payoff separated into two parts, the extrinsic payoff, which will still only depend on the distance between the strategies, and the intrinsic payoff which will depend only on the position of the player and, perhaps, the time. Consider, as in section \ref{SEC:Extension} a game with a discrete number of strategies where every pairwise interaction is subject to the same extrinsic payoff, given by the payoff matrix $A$ but there is an additional intrinsic payoff which is dependent on the player $x$, their strategy $u(x)$ and the time $t$. The payoff matrix therefore can be decomposed as
    
   \[A=B+C(x,t)\] 
    where $B$ is Toeplitz as before and $C(x,t)$ has constant rows for every $(x,t)\in \Omega_T$ (i.e. $C(x,t)=[c(x,t)\;c(x,t)\;\dots \;c(x,t)]$ where $c(x,t)$ is a column vector). We can rewrite our equation \eqref{ctsSpacePayoff} as
    \begin{equation*}
        w(x|u)=\int_\Omega K(x,y)u(x)^TBu(y)dy+\int_\Omega K(x,y)u(x)^T C(x,t)u(y)dy.
    \end{equation*}
    However, because of the form of $C(x,t)$,  $C(x,t)u(y)=c(x,t)\sum u(y)_i$.  In both the pure strategy and mixed strategy cases, $\sum u(y)_i=1$ so this is independent of $y$ and we can write 
    \begin{equation*}
        w(x|u)=\int_\Omega K(x,y)u(x)^TBu(y)dy+\langle u(x), c(x,t)\rangle \|K(x,\cdot)\|_{L^1(\Omega)}
    \end{equation*}

    In the same way as before, the extension into continuous strategy space is easy. the matrix $B$ and vector $c$ become infinite dimensional and we replace them with the functions $\rho$ and $F$ respectively. Thus we get our continuous strategy space functions
    \begin{equation}\label{inhompayoff}
        w(x|u)=\underbrace{\int_\Omega K(x,y)\rho(u(x)-u(y))dy}_{extrinsic}+ \underbrace{F(u(x,t),x,t)}_{intrinsic}
    \end{equation}

    We will now consider the requirements on $F$ that ensure that the model is still well founded. 

    \begin{proposition}\label{PROP:InhomogeneousModelWellFounded}
        Bounded strategy profiles of the game with players in $\Omega \subset \Rn$ choosing strategies in $\mathbb{R}$ by myopic best response with fitness as in \eqref{inhompayoff} will evolve as
        \begin{equation*}
            \frac{\partial}{\partial t}u(x,t)=g[u](x,t)+\partial_1 F(u(x,t),x,t)
        \end{equation*}
        Under the hypotheses of proposition \ref{PROP:ModelWellFounded} with the additional hypothesis that
        \begin{itemize}
            \item[(H4)] $F$ has a Lipschitz derivative with respect to its first argument and $F(z,x,t)\leq C_2z^2+A_3$ for some non-negative $C_2$ and $A_3$ uniformly in $\Omega_T$
        \end{itemize}
    \end{proposition}

    \begin{proof}
        As in proposition \ref{PROP:ModelWellFounded} we will consider a bounded strategy profile $u(\cdot, t):\Omega\to \mathbb{R}$. Every player will seek to update their strategy by some amount $h$ in order to take on their best response to $u(\cdot, t)$ after a time step of $\Delta t$ and in doing so they incur a cost of $\frac{h^2}{\Delta t}$. Let
        \begin{equation*}
            S_x(h):=\int_\Omega K(x,y)\rho(u(x,t)+h-u(y,t))dy+F(u(x,t)+h,x,t)
        \end{equation*}
        be the payoff player $x$ receives after changing their strategy by $h$. From the computations in the proof of proposition \ref{PROP:ModelWellFounded} we know that 
        \begin{equation*}
            \int_\Omega k(x,y)\rho(u(x,t)+h-u(y,t))dy\leq C_1h^2+B_1h+A_2
        \end{equation*}
        Additionally from $(H4)$ $F(u(x,t)+h,x,t)\leq {C}_2h^2+ A_3$ for some ${C}_2$ and ${A}_3.$ Therefore we have immediately that 
        
       \[S_x(h)\leq C_3h^2+A_4\] 

        We also need to show that $S_x(h)$ has a Lipschitz derivative. As in proposition \ref{PROP:ModelWellFounded} the first term of $S_x(h)$ has a locally Lipschitz derivative. Moreover, by $(H4)$ the second term also has a locally Lipschitz derivative with respect to $h$. 

        Having argued that $S_x(h)$ is uniformly subquadratic in $h$ and has a locally Lipschitz derivative in $h$ we note that the quantity $S_x(h)-\frac{h^2}{\Delta t}\leq C_4h^2+A_5$ for some negative $C_4$ when $\Delta t$ is sufficiently small. Therefore there is certainly a global maximizer $h^*$ and that maximizer satisfies 

       \[\frac{d}{dh}S_x(h^*)=2\frac{h^*}{\Delta t}.\]


        It is the same process as in proposition \ref{PROP:ModelWellFounded} that allows us to find that in the limit as $\Delta t\to 0$
        
       \[\frac{\partial}{\partial t}u(x,t)=\frac{1}{2}\frac{d}{dh}S_x(0)\]
        We note that
        
       \[\frac{\partial}{\partial h}S_x(0)=\int_\Omega K(x,y)\rho'(u(x,t)-u(y,t))dy+\partial_1F(u(x,t),x,t)\] 
       and do a trivial rescaling of space time to complete the proof.  
    \end{proof}

    For the proceeding let $f(u,x,t):=\partial_1F(u(x,t),x,t)$ and we are left with the simple inhomogeneous version of the nonlinear nonlocal diffusion problem we have been discussing
    \begin{equation}\label{inhom}
        u_t-g[u](x,t) =f(u,x,t)
    \end{equation}
    where $\rho$ and $f$ are both subquadratic for the game theoretic application to be sensible. It is not a difficult task to prove short time existence and uniqueness for solutions in this case. The theorem is written here and the proof is included in appendix \ref{APP:AdditionalInhomogeneousResults} because of its similarity to the proof of theorem \ref{THM:ShortTimeExistenceAndUniqueness}

    \begin{theorem}[Short time existence and Uniqueness for the inhomogeneous case]\label{THM:InhomogeneousExistenceAndUniqueness}
    The initial value problem $u_t-g[u]=f(u,x,t)$ has a unique continuous and bounded solution in $\Omega_T$ for some $T$ when $u(x,0)=u_0\in C_b^0(\Omega), \rho\in C^{1,1}(\mathbb{R})$, $K\in C^0_b(\Omega;L^1(\Omega))$ and $f$ is Lipschitz continuous with respect the the first variable with a Lipschitz constant which does not depend on time, is continuous with respect to space and time, and is bounded 
    \end{theorem}

    The proof proceeds practically identically to that of existence and uniqueness in the homogeneous case as it uses a BFPT argument. Finding global existence or even a finite time blow up result in the inhomogeneous case is certainly much harder. The Main issue is the nonlinearity of the nonlocality and so the standard Duhamel's property does not apply in this case. 

    \section{conclusion}\label{SEC:conclusion}
    In this study we have extended the game theoretic treatment of structured coordination to allow for continuous pure strategies and provide novel critical insights, extending the application areas beyond those for which only the traditional, discrete pure strategies of discrete player spaces are appropriate. With certain reasonable hypotheses on the components of the fitness function, we were able to show rigorously that through a myopic best response like update rule the situation can be modeled in continuous time by way of a nonlinear nonlocal equation, similar to existing nonlocal diffusion problems. 

    The nonlinearity in this model prevents the use of Fourier analysis, semi-group theory, or comparison principles in proving our results, but through elementary analysis and PDE techniques we were able to determine short time existence and uniqueness for the general setting. With some additional requirements on $\rho$, we can strengthen this to find global existence and uniqueness. In addition to these results, we found that solutions with Lipschitz initial data remain Lipschitz continuous although the modulus of continuity may increase exponentially. In the special case of the Cauchy problem, the Lipschitz constant does not depend on the shape of the kernel, $K$ as long as it is translation invariant. 

    After giving several trivial analytical examples, we showed that simple numerical methods are stable and convergent, which allows us to visually examine solutions to the initial value problem without boundary data prescribed. Using these results we were able to carry out several numerical examples which supported the analytical results about stationary solutions. Finally, we considered the inhomogeneous problem and again showed short time existence and uniqueness of solutions to the IVP with no boundary data.

    Not only do these results help us understand the model as a way of discussing coordination in space, but they also represent advances in our understanding of nonlinear, nonlocal diffusion problems. On the modeling side, we see examples of solutions seeming to converge towards non-constant equilibria, the continuous analog of the non-trivial equilibria in the discrete case. Moreover, when discontinuities in strategy emerge, we can determine how quickly and how severely they can appear.  By characterizing stationary solutions, we can significantly advance our understanding of both the classical game and nonlinear nonlocal diffusion problems. The concentration of strategies appears as a smoothening in some parts of the domain and tends to create discontinuities in other parts of the domain. Further study into nonlinear nonlocal diffusion problems is required to describe the asymptotic behavior of solutions and the stability of stationary solutions. From this study, it is clear that understanding coordination as a nonlinear nonlocal diffusion problem allows us to examine the system in exciting and novel ways making accessible insights that were previously impossible.  



Both J.S.M. and T.A.M. contributed to the formal analysis and investigation. J.S.M. and N.H.F. Conceptualized the work. J.S.M. wrote the software and wrote the original draft.\newline

The authors declare no competing financial interests.

\section{\label{DataAvail}Data Availability Statement}
All of the software which was used to run the numerical simulations and generate the figures in this manuscript can be found at the repository linked here: \url{https://github.com/feffermanlab/JSM_2024_ContinuousCoordination/releases/tag/v1.0.0}. It is also cited as \cite{McAlister2025}.

\appendix  
\section{Alternative Proof for Global Existence}\label{APP:AltProofForGlobalEx}
Recall theorem \ref{THM:GlobalExistenceAndUniqueness} when $\rho\in C^{1,1}$ satisfies the coordination property \eqref{coordinationcondition}. Under these conditions the a solution to the the initial value problem $u_t=g[u]$ exists and is unique for all finite time. The short proof is included in the main text, but there is also an alternative proof which reveals the repeatability of the extension principle as described in the proof of theorem \ref{THM:ShortTimeExistenceAndUniqueness}

    \begin{theorem*}[Global existence and uniqueness with particular $\rho\in C^{1,1}$]
        Let $\rho\in C^{1,1}(\mathbb{R})$ satisfy \eqref{coordinationcondition}. Under this strengthened hypothesis, the Initial Value Problem $u_t=g[u]$ with $u(x,0)=u_0\in C_b^0(\Omega)$ has a unique continuous and bounded solution for all finite time.
    \end{theorem*}

    \begin{proof}
        We modify our proof from Theorem \ref{THM:ShortTimeExistenceAndUniqueness}. Equip the function space $C_b^0(\overline{\Omega}_T)$ with the standard sup norm $\|u\|=\sup_{t\in [0,T]}\|u(\cdot ,t)\|_\infty$. Now let \[E_{T}:=\{u\in C_b^0(\overline{\Omega}_T;\mathbb{R});u(x,0)=u_0,\|u\|\leq \|u_0\|\}\] for some $T$ to be determined later. Observe again that $u(x,t)\equiv u_0(x)$ is in $E_{T}$ so it is non-empty. Also, observe that $E_{T}$ is complete. Consider the same operator $\Theta:C_b^0(\overline{\Omega}_T,\mathbb{R})\rightarrow C_b^0(\overline{\Omega}_T,\mathbb{R})$ as in theorem \ref{THM:ShortTimeExistenceAndUniqueness}. We will show again that $\Theta:E_{T}\rightarrow E_{T}$.

        We again use lemma \ref{LEM:welldefined} to say that $\Theta u$ is clearly continuous and indeed continuously differentiable in time. 
        To say that $\|\Theta u\|\leq \|u_0\|_\infty$ we need only repeat our argument from lemma \ref{LEM:MaximumPrinciple}.

        Suppose there is a time $t_0$ and position $x_0$ where $\Theta u(x_0,t_0)-\epsilon t_0\geq  v:=\|u_0\|+\epsilon$. $\Theta u$ is differentiable in time so, supposing that $t_0$ is the first time this inequality is satisfied we can say that $\frac{\partial}{\partial t}\Theta u(x_0,t_0))-\epsilon\geq 0$. We also have that $\frac{\partial}{\partial t}(\Theta u (x_0,t_0))=g[u](x_0,t_0)$. We already know that $u$ and $\Theta u$ are continuous so $u(x,t_0)\leq v$ for all $x\in \mathbb{R}^n$. Because $\rho'(z)\leq 0$ whenever $z\geq 0$ we know that $\rho'(u(x_0,t_0)-u(y,t_0))\leq 0$ for all $y$ and thus $g[u](x_0,t_0)\leq 0$. Therefore we get a contradiction 
        \begin{equation*}
            0\geq g[u](x_0,t_0)=\frac{\partial}{\partial t}(\Theta u (x_0,t_0))\geq \epsilon >0
        \end{equation*}
        Thus we conclude that $\|\Theta u(\cdot,t)-\epsilon t\|_\infty\leq \|u_0\|+\epsilon$ for all $t\in [0,T)$.  This inequality holds for any epsilon so we have shown that regardless of the choice of $T$, $\|\Theta u\|\leq \|u_0\|_\infty$.
        Therefore we have that $\Theta: E_{T}\rightarrow E_{T}$ 

        Next we must show that $\Theta$ is a contraction on $E_{T}$. The argument here is exactly the same as the argument in theorem \ref{THM:ShortTimeExistenceAndUniqueness}. We will have the Lipschitz constant for $g$, $2L_{\|u_0\|}$ by lemma \ref{LEM:Lipschitzg} so we get that
        
       \[\|\Theta u-\Theta v\|\leq C^gT\|u-v\|.\] We are assured that such a $C^g$ exists because we are working in a compact subset of $C^0_b(\Omega_T)$, namely $E_T$.
        As before we now know that if $T\leq \frac{1}{2C^g}$ then $\Theta$ is a contraction from $E_{T}$ to $E_T$ and therefore there is a unique solution to the IVP.

        Now we want to show that we can extend this solution to any finite time. 
        consider any $u_0\in C^0_b(\Omega)$ and note that we know that a solution exists and is unique on $[0,T)$. Take $u(\cdot, T-\eta)$ for some $\eta>0$ to be our new initial condition and note again that because $\|u(\cdot, T-\eta)\|\leq \|u_0\|$ the original Lipschitz constant for $g$ given the bound from $u_0$, is still an appropriate Lipschitz constant for $g$ given the bound on the new initial data $\|u(\cdot,T-\eta)\|$. Therefore we can prove existence for the same length of time, $T=\frac{1}{2C^g}$, and we have a solution on $[T-\eta, 2T-\eta)$. Moreover, when the solutions overlap they are identical so the solution on $[0,T)$, overlaps perfectly with the solution on $[T-\eta, 2T-\eta)$ and is continuous and differentiable in time, therefore, it is a solution on $[0,2T-\eta)$. We can repeat this process any number of times to show that our solution exists for all finite times. This completes the proof of global existence and uniqueness. 
    \end{proof}

\section{Additional Numerical Result}\label{APP:AdditionalNumericalResults}
    \begin{lemma}[discrete maximum principle (Forward Euler)]\label{LEM:DiscreteMaximumPrinciple}
        If $\Omega_T$ is a bounded time cylinder with discretization $\Omega_T^{(h,\tau)}$, and $w\in \mathcal{V}(\Omega_T^{(h,\tau)})$ which satisfies \eqref{ForwardEuler} with $K(x,y)\in C^0_b(\Omega;C^0_b)$ with $w(\mathbf{x},0)=\pi^h u_0(\mathbf{x})$, then, when $\tau\leq L_\rho\sup_{\mathbf{x}\in \Omega^h}\|\pi^hK(\mathbf{x},\cdot)\|_{\ell^\infty(\Omega)}$ 
       \[\|w\|_{\ell^\infty\left({\Omega^{(h,\tau)}_T}\right)}\leq \|u_0\|_{L^\infty(\Omega)}\]
        where $L_\rho$ is the Lipshitz constant for $\rho'$ on $[-2\|u_0\|_{L^\infty(\Omega)},2\|u_0\|_{L^\infty(\Omega)}]$. 
    \end{lemma}

    \begin{proof}
        Suppose that $\|w(\cdot, t_i)\|_{\ell^\infty}= M$. We seek to show that $w(\mathbf{x},t_{i+1})\leq M$ for all $\mathbf{x}\in \Omega^h$.  Consider $\mathbf{x}\in \Omega^h$, we will break the proof of this claim into two cases, the first if $w(\mathbf{x},t_i)=M$ and the second, $w(\mathbf{x},t_i)<M$. 

        In the first case, observe that, because $w(\mathbf{x},t_i)-w(\mathbf{y},t_i)\geq 0$ for all $\mathbf{y}\in ^-\Omega^h$, we have that $\rho'(w(\mathbf{x},t_i)-w(\mathbf{y},t_i))\leq 0$ for all $\mathbf{y}\in ^-\Omega^h$. It is clear then that
        
       \[w(\mathbf{x},t_{i+1})=M+\tau h^n\sum_{\mathbf{y}\in ^-\Omega^h}K(\mathbf{x},\mathbf{y})\rho'(u(\mathbf{x},t_i)-u(\mathbf{y},t_i))\leq M\]

        Now we consider the case that $w(\mathbf{x},t_{i})<M$. If we show that 
        \begin{equation}\label{DMPcase2} \tau h^n\sum_{\mathbf{y}\in ^-\Omega^h}K_\epsilon(\mathbf{x}-\mathbf{y})\rho'(u(\mathbf{x},t_i)-u(\mathbf{y},t_i))\leq M - w(\mathbf{x},t_i)\end{equation}
        then surely $w(\mathbf{x},t_{i+1})\leq M$. We can show this by observing that $\rho'$ is Lipschitz with Lipschitz constant $L_\rho$ on the range of $w$ and that, $w(\mathbf{x},t_i)-w(\mathbf{y},t_i)\geq w(\mathbf{x},t_i)-M.$ Because of the assumption that $\rho'(z)\geq 0$ when $z<0$ and $\rho'(z)\leq 0$ when $z>0$, in order to show the upper bound of the sum in \eqref{DMPcase2} we assume the worst case which is that $w(\mathbf{x},t_i)-w(\mathbf{y},t_i))\leq 0$ for all $\mathbf{y}\in ^-\Omega^h$. In this worst case we know that
        
       \[0\leq \rho'(w(\mathbf{x},t_i)-w(\mathbf{y},t_i))\leq L_\rho(M-w(\mathbf{x},t_i))\] for each $\mathbf{y}\in ^-\Omega^h$. Thus we have that
        \begin{equation*}
        \begin{split}
            \tau h^n\sum_{\mathbf{y}\in ^-\Omega^h}&K(\mathbf{x},\mathbf{y})\rho'(u(\mathbf{x},t_i)-u(\mathbf{y},t_i))\\
            &\leq \tau h^nL_\rho(M-w(\mathbf{x},t_i))\sum_{\mathbf{y}\in ^-\Omega^h}K(\mathbf{x},\mathbf{y})\\
            &\leq \tau L_\rho D(M-w(\mathbf{x},t_i))
            \end{split}
        \end{equation*}
        Where $D=\sup_{\mathbf{x}\in \Omega^h}\|\pi^hK(\mathbf{x},\cdot)\|_{\ell^\infty(\Omega)}$. Therefore, when $\tau<\frac{1}{L_\rho D}$, inequality \eqref{DMPcase2} is satisfied, so we have proved that, whenever $\tau <\frac{1}{L\rho D}$, $\|w(\cdot, t_{i+1})\|_{\ell^\infty(\Omega^h)}\leq \|w(\cdot, t_{i})\|_{\ell^\infty(\Omega^h)}$. The desired result is an obvious consequence. 
        
    \end{proof}



\section{Additional Inhomogeneous Result}\label{APP:AdditionalInhomogeneousResults}
\begin{theorem}[Short time existence and uniqueness for the inhomogeneous problem]The initial value problem $u_t-g[u]=f(u,x,t)$ has a unique continuous and bounded solution in $\Omega_T$ for some $T$ when $u(x,0)=u_0\in C_b^0(\Omega), \rho\in C^{1,1}(\mathbb{R})$, $K\in C^0_b(\Omega;L^1(\Omega))$ and $f:\mathbb{R}\times \Omega\times \mathbb{R}\to \mathbb{R}$ is Lipschitz continuous with respect the the first variable with a Lipschitz constant, $L_f$, which does not depend on time, is merely continuous in space and time, and is bounded. 
\end{theorem}
\begin{proof}
    Let $\Omega_T=\Omega\times [0,T]$ with $T$ to be chosen later. Equip the function space $C^0_b(\Omega_T)$ with the standard sup norm $\|u\|=\sup_{t\in [0,t]}\|u(\cdot, t)\|_\infty.$ Now for some $R>\|u_0\|_\infty$ let $E_{R,T}:=\{u\in C^0_b(\Omega_T,\mathbb{R});u(x,0)=u_0,\|u\|\leq R\}$ and observe that $u(x,t)=u_0(x)$ for all $t$ is in this closed subset of the Banach space. As before we will use the BFPT. Clearly a solution to the IVP will satisfy
    \begin{equation}\label{inhomIntegralEquation}
        u(x,t)=u_0(x)+\int_0^tg[u](x,s)+f(u(s),x,s) ds
    \end{equation}
    Let $\Psi:C^0_b(\Omega_T,\mathbb{R})\to C^0_b(\Omega_T,\mathbb{R})$ where $\Psi u = u_0+\int_0^tg[u](x,s)+f(u(s),x,s)ds$. Clearly if there is a solution to $\Psi u=u$ then we have a solution to \eqref{inhomIntegralEquation} and thus the IVP.

    We start by showing that $\Psi:E_{R,T}\to E_{R,T}$. By lemma \ref{LEM:welldefined} we know that $g[u]$ is continuous so its time antiderivative is clearly continuous. Moreover, by assumption $f$ is continuous so its time antiderivative will also be continuous.

    To show that $\|\Psi u\|\leq R$  we note that both $g[u]$ and $f$ are bounded. $g$ is bounded exactly as described in theorem \ref{THM:ShortTimeExistenceAndUniqueness} so there is a $T_R$ before which 
    \[\sup_{t\in[0,T)}\|\int_0^tg[u](\cdot,s)ds\|_\infty \leq \frac{R-\|u_0\|_\infty}{2}.\] 
    $f(u(x,s),x,s)$ is bounded, by assumption, so likewise there is a $T_{Rf}$ so that \[\int_0^tf(u(x,s),x,s)ds<\frac{R-\|u_0\|_\infty}{2}\] when $t<T_{Rf}$. Thus there is a $T_1$ so that when $T<T_1$, $\Psi:E_{R,T}\to E_{R,T}$. 

    To show that this is a contraction consider 
    \begin{equation*}
        \begin{split}
            \|\Psi u-\Psi v\|&=\sup_{t\in[0,T)}\bigg\|\int_0^tg[u](\cdot,s)-g[v](\cdot,s)ds+\int_0^tf(u(\cdot,s),\cdot,s)-f(v(\cdot,s),\cdot,s)ds\bigg\|_\infty\\
            &\leq \int_0^T\|g[u](\cdot,s)-g[v](\cdot,s)\|_\infty ds+\int_0^T|f(u(\cdot,s),\cdot,s)-f(v(\cdot,s),\cdot,s)|ds\\
            &\leq TC^g\|u-v\|+L_fT\|u-v\|
        \end{split}
    \end{equation*}
    Thus if $T<\frac{1}{2(c^g+L_f)}$ then $\|\Psi u -\Psi v\|\leq \frac{1}{2}\|u-v\|$. 

    Therefore when $T<\min\{T_1,\frac{1}{2(c^g+L_f)}\}$ then $\Psi:E_{R,T}\to E_{R,T}$ is a contraction and thus it has a unique fixed point. Thus we have shown that there is a unique solution to the integral equation \eqref{inhomIntegralEquation} and thus the IVP for some short time. We can extend this in the same way as in \ref{THM:ShortTimeExistenceAndUniqueness} but can not extend this to global existence without a maximum principle.
    \end{proof}


% If in two-column mode, this environment will change to single-column
% format so that long equations can be displayed. Use
% sparingly.
%\begin{widetext}
% put long equation here
%\end{widetext}



% Create the reference section using BibTeX:
\bibliography{ContinuousCoordination}

\end{document}
%
% ****** End of file apstemplate.tex ******

